\chapter{The Source Stage 1 Route}
\label{ch:stage1-source}

\section{The source objects}

\begin{definition}[Initial theta and theater data]
  \label{def:initial-theta-data}
  \lean{Iut.InitialThetaData}
  \lean{Iut.ThetaLTorsionRepresentationData.IsKernelField}
  \lean{Iut.ThetaLTorsionRepresentationData.ImageContainsSL2}
  \lean{Iut.ThetaQParameterOrderData.OrdersPrimeToL}
  \leanok
  The source construction begins with the initial \(\Theta\)-data of \IUT{} I:
  a number field, an elliptic curve, bad places, label data, and the
  arithmetic hypotheses from which Hodge theaters are built.  The mod-\(l\)
  representation, its kernel field, containment of
  \(\mathrm{SL}_2(\mathbb Z/l\mathbb Z)\), and positive bad-place
  q-parameter orders prime to \(l\) are explicit predicates and data rather
  than free propositions.
\end{definition}

\begin{definition}[Hodge theater histories]
  \label{def:hodge-theater-histories}
  \lean{Iut.ThetaArithmeticHistory}
  \lean{Iut.Stage1HodgeTheaterSource}
  \lean{Iut.Stage1HodgeTheaterSource.ofInitialThetaData}
  \lean{Iut.Stage1HodgeTheaterSource.histories_not_identified}
  \leanok
  A Hodge theater carries an arithmetic holomorphic structure.  The Stage 1
  source now constructs the domain history of the \qPilot{} and the
  codomain history of the \ThetaPilot{} as distinct constructors.  Their
  separation is therefore a theorem of the datatype, not an inequality between
  caller-supplied strings.
\end{definition}

\begin{definition}[Prime strips and links]
  \label{def:prime-strip-links}
  \lean{Iut.ThetaPrimeStrip}
  \lean{Iut.ThetaBridgeSource}
  \lean{Iut.ThetaPlusMinusEllNFHodgeTheater}
  \lean{Iut.Stage1HodgeTheaterSource.zeroColumn_primeStrip_compatible}
  \lean{Iut.Stage1HodgeTheaterSource.oneColumn_primeStrip_compatible}
  \leanok
  The \(\Theta^{\times\mu}_{\mathrm{LGP}}\)-link between the \(0\)- and
  \(1\)-columns has a value-group portion and a unit portion.  The value-group
  portion relates \qPilot{} data at \((0,0)\) to \ThetaPilot{} data at
  \((1,0)\).  At the IUT I boundary, prime strips are indexed by their owner
  history and project the global orbicurve, selected places, local orbicurves,
  cusps, and q-orders from the same initial theta datum.  The
  \(\Theta^{\pm\ell}\) and \(\Theta^{\mathrm{NF}}\) components are glued by an
  equality of their typed gluing-role prime strips.  The later \(\IPL\)
  property remains part of the Theorem 3.11 route.
\end{definition}

\begin{definition}[Topological pseudo-monoids]
  \label{def:source-topological-pseudo-monoids}
  \lean{Iut.SourceTopologicalPseudoMonoid}
  \lean{Iut.SourceTopologicalPseudoMonoid.MultiplicationDomain}
  \lean{Iut.SourceTopologicalPseudoMonoid.partialMul}
  \lean{Iut.SourceTopologicalPseudoMonoid.continuous_partialMul}
  \lean{Iut.SourceTopologicalPseudoMonoid.partialMul_assoc}
  \lean{Iut.SourceTopologicalPseudoMonoid.IsDivisible}
  \lean{Iut.SourceTopologicalPseudoMonoid.IsCyclotomic}
  \lean{Iut.SourceTopologicalPseudoMonoid.Iso.maps_domain_iff}
  \lean{Iut.SourceTopologicalGroupPseudoMonoidActionPair}
  \lean{Iut.SourceTopologicalGroupPseudoMonoidActionPair.action_maps_domain_iff}
  \lean{Iut.SourceTopologicalGroupPseudoMonoidActionPair.action_partialMul}
  \lean{Iut.SourceTopologicalGroupPseudoMonoidActionPair.fixedPseudoMonoid.mul_mem_iff}
  \lean{Iut.SourceProfinitePseudoMonoidActionPair.quotientAction_comp}
  \lean{Iut.SourceProfinitePseudoMonoidActionPair.descend}
  \lean{Iut.SourceProfinitePseudoMonoidActionPair.comap_descend}
  \leanok
  Following IUT~I, Section~0, a topological pseudo-monoid is the embedded
  image of a topological space in a topological abelian group.  Its product is
  defined exactly for pairs whose ambient product remains in that image; the
  image need not contain the identity and need not be closed under arbitrary
  products.  The restricted product is continuous, commutative, and
  associative whenever both composites are defined.  Divisibility retains
  both ambient roots and the power-membership equivalence, while
  cyclotomicity retains the torsion identification with \(\mathbb Q/\mathbb Z\)
  and stability under torsion.  Continuous morphisms preserve defined
  products, isomorphisms reflect their domain, and continuous group actions
  are by these pseudo-monoid automorphisms.  For Definition~5.2(v), the
  invariant carrier is constructed as a sub-pseudo-monoid and is proved to
  inherit every product defined in the original carrier.  A profinite action
  whose kernel acts trivially descends through a surjective profinite
  quotient; continuity follows from the compact-to-Hausdorff quotient
  topology, and pullback of the descended action recovers the original one.
  These results implement the fixed-point and
  $\pi_1^{\kappa\text{-sol}}$ factorization mechanisms.  They do not construct
  the ambient overgroup, rational fundamental group, or the coric
  rational-function carriers supplied by the cited absolute-anabelian
  reconstruction algorithms.
\end{definition}

\begin{definition}[Source Frobenioid prime strips]
  \label{def:source-frobenioid-prime-strips}
  \lean{Iut.PreFrobenioid.RationalMonoidTransport}
  \lean{Iut.PreFrobenioid.CharacteristicSplitting}
  \lean{Iut.SplitFrobenioidPresentation}
  \lean{Iut.SplitTopologicalMonoidPresentation}
  \lean{Iut.SourceFPrimeStrip}
  \lean{Iut.IUTINonarchimedeanLocalModel.dToCoreOpenEmbedding}
  \lean{Iut.IUTINonarchimedeanLocalModel.dToCore_injective}
  \lean{Iut.IUTINonarchimedeanLocalModel.dToCore_open_range}
  \lean{Iut.SourceFinitePlaceReconstruction.carrierEquivAlgebraicClosure}
  \lean{Iut.SourceFinitePlaceReconstruction.continuous_stageIntegralToInd}
  \lean{Iut.SourceFinitePlaceReconstruction.exists_stage_representation}
  \lean{Iut.SourceFinitePlaceReconstruction.indIntegralMonoidEquivSourceMLF}
  \lean{Iut.SourceFinitePlaceReconstruction.indIntegralMonoidEquivSourceMLF_galois}
  \lean{Iut.SourceFinitePlaceReconstruction.carrierMulEquiv_galoisContinuousLinearMap}
  \lean{Iut.SourceFinitePlaceReconstruction.continuous_indGaloisMulAut}
  \lean{Iut.SourceFinitePlaceReconstruction.indGaloisContinuousMulEquiv_mul}
  \lean{Iut.SourceFinitePlaceReconstruction.continuous_finiteStageGaloisEvaluation}
  \lean{Iut.SourceFinitePlaceReconstruction.continuous_stageGaloisAction}
  \lean{Iut.SourceFinitePlaceReconstruction.continuous_galoisLinearMap_comp_stageToCarrier}
  \lean{Iut.SourceFinitePlaceReconstruction.continuous_indGaloisAction_comp_stageIntegralToInd}
  \lean{Iut.SourceFinitePlaceReconstruction.stageIntegralTransition_trans}
  \lean{Iut.SourceFinitePlaceReconstruction.continuous_stageGaloisIntegralAction}
  \lean{Iut.SourceFinitePlaceReconstruction.stageIntegralTransition_galois}
  \lean{Iut.SourceFinitePlaceReconstruction.integralIndSystem}
  \lean{Iut.SourceFinitePlaceReconstruction.stageIntegralToIndHom_transition}
  \lean{Iut.SourceFinitePlaceReconstruction.stageIntegralToIndHom_galois}
  \lean{Iut.SourceFinitePlaceReconstruction.integralIndSystemLimit}
  \lean{Iut.SourceFinitePlaceReconstruction.integralIndSequentialPresentation}
  \lean{Iut.SourceFinitePlaceReconstruction.denseRange_algebraMap_base}
  \lean{Iut.SourceFinitePlaceReconstruction.exists_stage_primitive_generator}
  \lean{Iut.SourceFinitePlaceReconstruction.exists_global_polynomial_root_same_adjoin}
  \lean{Iut.SourceFinitePlaceReconstruction.exists_global_polynomial_root_generating_stage}
  \lean{Iut.SourceFinitePlaceReconstruction.globalPolynomialRootStage_surjective}
  \lean{Iut.SourceFinitePlaceReconstruction.stageIndexCountable}
  \lean{Iut.SourceFinitePlaceReconstruction.integralIndSourceSequentialPresentation}
  \lean{Iut.SourceFMonoAnalyticPrimeStrip}
  \lean{Iut.SourceFGloballyRealifiedMonoAnalyticPrimeStrip}
  \lean{Iut.sourceFPrimeStripToD}
  \lean{Iut.SourceFPrimeStripCapsule}
  \lean{Iut.SourceThetaHodgeTheater.associatedF}
  \lean{Iut.SourceThetaHodgeTheater.MonoAnalyticTransport}
  \lean{Iut.SourceThetaHodgeTheater.associatedFMonoAnalytic}
  \lean{Iut.SourceFThetaBridge}
  \lean{Iut.SourceFThetaBridge.associatedD}
  \lean{Iut.SourceFThetaBridge.associatedD_lift_unique}
  \leanok
  The transitive Frobenioids I dependency is retained explicitly:
  \(O^\triangleright(-)\) is a contravariant functor on the isotropic
  linear category, and a characteristic splitting is a subfunctor whose
  product with the units is objectwise bijective and satisfies the
  isotropic-hull condition.  Thus a split Frobenioid is not identified with
  a Frobenius-degree trivialization.  At nonarchimedean places an
  \(F^{\sim}\)-component is an actual split Frobenioid; at archimedean
  places it also carries the compact/noncompact \(TM^{\sim}\)
  factorization, tied to the selected characteristic splitting.
  Example~3.2(i) identifies the group represented by \(D_v\) with the
  arithmetic fundamental group of \(X_v\).  The source model now uses exactly
  this group, and the sign-quotient geometry constructs its open embedding
  into the arithmetic fundamental group of \(C_v\), with injectivity and open
  image proved.  This is the geometric target of Definition~5.2(v), not yet
  the absolute-anabelian algorithm recovering it from an abstract \(D_v\)
  group.
  For an actual finite place, the algebraic closure is now the filtered
  topological-module colimit of its finite Galois local subfields.  Restricting
  this colimit to its nonzero integral elements gives continuous finite-stage
  inclusions that exhaust the carrier, a multiplicative equivalence with
  \(O_{\overline{K_v}}^\triangleright\), and the literal compatible absolute
  Galois action.  Normality of each finite stage now gives a continuous
  restriction of every Galois automorphism; the colimit universal property
  produces the corresponding continuous ind-automorphism, and Lean proves
  that these automorphisms satisfy the identity and composition laws and
  restrict continuously in both directions to
  \(O_{\overline{K_v}}^\triangleright\).  In accordance with the meaning of
  the prefix ``ind-'' in Absolute Anabelian Topics~III, Section~0, Lean also
  retains the finite Galois integral stages as an actual filtered system of
  topological monoids.  Its transitions are continuous and equivariant;
  continuity of Krull restriction to the finite Galois group, together with
  its discrete topology, proves joint continuity of evaluation on every
  stage.  The stage inclusions are injective, exhaust
  \(O_{\overline{K_v}}^\triangleright\), and commute with the limit action.
  The global number field is dense in \(K_v\).  Primitive elements,
  same-degree polynomial approximation, continuity of roots, and Krasner's
  lemma show that every finite Galois stage is generated by a root of a
  polynomial over the global field.  These polynomial-root witnesses form a
  countable type and surject onto the stage type.  The countable
  filtered-category theorem therefore constructs a monotone cofinal
  natural-number subsystem, hence the strict positive-integer presentation
  used by the paper.  Reconstruction from the abstract \(D_v\) group remains
  open.
  A globally realified strip additionally carries the global Frobenioid,
  its prime/place bijection, all local characteristic and realified monoids,
  and commuting currency-exchange isomorphisms \(\rho_v\), compared
  componentwise with the Example~3.5 model.
  Holomorphic \(F\)-prime-strips form the isomorphism-only category required
  by Definition~5.2(iii), hence admit genuine capsules, and their associated
  \(D\)-prime-strips are functorial.  The \(F\)-strip associated to a
  Definition~3.6 theta-Hodge theater is derived, not supplied independently.
  Its mono-analytic \(F^\vdash\)-strip is constructed only from an explicit
  structure-preserving transport of the rational monoids, characteristic
  splittings, and selected archimedean splitting; equivalences of the bare
  carrier categories are insufficient.
  Consequently, selected \(F\)-level theta-bridge arrows induce a
  \(D\)-theta bridge, with the Corollary~5.3 uniqueness of selected lifts
  exposed as a separate theorem.  Proposition~2.2's construction of the
  rational-monoid transport, the Examples~3.2--3.5 model constructions,
  mono-analyticization, global realification, the remaining reconstruction
  content of Definition~5.2(v), and clauses~(vi)--(viii) remain explicit
  source obligations.
\end{definition}

\begin{definition}[Capsules and processions]
  \label{def:source-processions}
  \lean{Iut.CategoryCapsule}
  \lean{Iut.CategoryCapsule.FullPolyMorphism}
  \lean{Iut.CategoryProcession}
  \lean{Iut.CategoryProcession.Hom}
  \lean{Iut.CategoryProcessionObject}
  \lean{Iut.IUTIPositiveLabel.signLabelEquiv}
  \lean{Iut.SourceDThetaBridgeCore}
  \lean{Iut.SourceDThetaBridgeCore.procession}
  \lean{Iut.SourceDThetaBridgeCore.procession_label_indeterminacy}
  \lean{Iut.SourceDMonoAnalyticization}
  \lean{Iut.SourceDThetaBridgeCore.monoAnalyticProcession}
  \lean{Iut.SourceDThetaBridgeCore.monoAnalyticProcession_label_indeterminacy}
  \leanok
  The capsule conventions of \IUT{} I, Section 0 and Definition~4.10 are
  formalized directly.  A capsule-full poly-morphism contains all
  componentwise isomorphisms over a fixed injection.  An \(n\)-procession has
  a \(j\)-capsule in position \(j\), and its morphisms use order-preserving
  injections together with capsule-full poly-morphisms.  Capsules and
  processions carry the identity and composition operations required by the
  categories used in Proposition~4.11.  The positive representatives
  \(1,\ldots,l^\ast\) are proved equivalent to the nonzero Tate labels modulo
  sign.  Consequently, a Proposition~4.7 evaluation-label bijection derives
  the nested holomorphic \(l^\ast\)-procession and its exact \(l^\ast!\)
  indeterminacy.  An actual natural mono-analyticization functor maps this to
  the mono-analytic procession while preserving the cardinality and factorial
  proofs.  Construction of that functor from the source subcategories,
  derivation of the label bijection from evaluation sections, and bridge
  functoriality remain source obligations.
\end{definition}

\begin{definition}[Toy degree-preorder realization]
  \label{def:iut-i-theater-realization}
  \uses{def:initial-theta-data,def:hodge-theater-histories,def:prime-strip-links}
  \lean{Iut.ToyIUTIHodgeTheaterRealization}
  \lean{Iut.ToyIUTIHodgeTheaterRealization.ofStage1HodgeTheaterSource}
  \lean{Iut.Stage1HodgeTheaterSource.hasToyIUTIRealization}
  \lean{Iut.Stage1.IUTStage1InitialThetaPrimitiveSourcePacket}
  \leanok
  The typed source constructs local and global effective-divisor
  degree preorder categories, the theta-link equivalence, and the
  \(\Theta^{\pm\ell}/\Theta^{\mathrm{NF}}\) gluing equivalences.  This
  degree-level test double does not prove Definition~3.6, Corollary~3.7,
  Definitions~4.1 and~5.2, Proposition~6.7, Remark~6.12.2, or
  Definition~6.13 of \IUT{} I.  Frobenioid and Kummer structure is not
  asserted by this model.
\end{definition}

\begin{definition}[Tempered semi-graph cusp rigidity]
  \label{def:tempered-semigraph-cusp-rigidity}
  \lean{Iut.SourceSemiGraph}
  \lean{Iut.SourceSemiGraph.verticialCardinality}
  \lean{Iut.SourceSemiGraph.isClosed_iff_verticialCardinality_eq_two}
  \lean{Iut.SourceSemiGraph.Hom}
  \lean{Iut.SourceSemiGraph.Hom.incidentBranchMap}
  \lean{Iut.SourceSemiGraph.Hom.IsImmersion}
  \lean{Iut.SourceSemiGraph.Hom.IsExcision}
  \lean{Iut.SourceSemiGraph.Hom.IsProper}
  \lean{Iut.SourceSemiGraph.Hom.isProper_iff_preserves_verticialCardinality}
  \lean{Iut.SourceSemiGraph.Hom.IsGraphCovering}
  \lean{Iut.SourceSemiGraph.Sub}
  \lean{Iut.SourceSemiGraph.compactification}
  \lean{Iut.SourceSemiGraph.compactification_isGraph}
  \lean{Iut.SourceSemiGraph.Hom.compactificationMap}
  \lean{Iut.SourceSemiGraph.Action}
  \lean{Iut.SourceSemiGraph.Action.branchStabilizer_le_edgeStabilizer}
  \lean{Iut.SourceSemiGraph.Action.branchStabilizer_le_vertexStabilizer}
  \lean{Iut.SourceSemiGraph.Action.actIncidentBranch_bijective}
  \lean{Iut.SourceCofilteredFiniteActionSystem}
  \lean{Iut.SourceCofilteredFiniteActionSystem.actSection}
  \lean{Iut.SourceCofilteredFiniteActionSystem.sectionAction}
  \lean{Iut.SourceCofilteredFiniteActionSystem.sectionStabilizer}
  \lean{Iut.SourceCofilteredFiniteActionSystem.mem_sectionStabilizer_iff_forall_coordinate}
  \lean{Iut.SourceCofilteredFiniteActionSystem.sectionStabilizer_eq_iInf_coordinateStabilizer}
  \lean{Iut.SourcePointedAnabelioidHom}
  \lean{Iut.SourcePointedAnabelioidHom.fundamentalGroupHom}
  \lean{Iut.SourcePointedAnabelioidHom.fundamentalGroupHom_comp}
  \lean{Iut.SourcePointedAnabelioidHom.TwoIso}
  \lean{Iut.SourcePointedAnabelioidHom.TwoIso.fundamentalGroupHom_eq_conjugate}
  \lean{Iut.SourcePointedAnabelioidHom.fundamentalGroupImageContinuousMulEquiv}
  \lean{Iut.continuousAction_preGaloisCategory}
  \lean{Iut.continuousAction_fiberFunctor}
  \lean{Iut.continuousAction_isFundamentalGroup}
  \lean{Iut.continuousActionEtaleFundamentalGroup}
  \lean{Iut.continuousActionPointedHom}
  \lean{Iut.continuousActionPointedHom_fundamentalGroupHom}
  \lean{Iut.SourceConnectedFiniteEtaleHom}
  \lean{Iut.SourceConnectedFiniteEtaleHom.ofProfiniteOpenEmbedding}
  \lean{Iut.SourceConnectedFiniteEtaleHom.ofProfiniteOpenEmbedding_fundamentalGroupHom}
  \lean{Iut.SourcePaperInducedSet.equivStandard}
  \lean{Iut.SourcePaperInducedSet.toStandard_smul}
  \lean{Iut.sourceInducedSet_stabilizer_isOpen}
  \lean{Iut.sourceInducedAction}
  \lean{Iut.sourcePaperInducedAction}
  \lean{Iut.sourcePaperInducedActionIso}
  \lean{Iut.sourceOpenCosetSliceEquivalence}
  \lean{Iut.sourceSliceForgetAdjProduct}
  \lean{Iut.sourceInductionRestrictionAdjunction}
  \lean{Iut.sourceOpenSubgroupRestrictionSliceIso}
  \lean{Iut.SourceFiniteEtaleFunctorFactorization}
  \lean{Iut.sourceOpenSubgroupFiniteEtaleFactorization}
  \lean{Iut.SourcePointedSemiGraphOfAnabelioids}
  \lean{Iut.SourcePointedSemiGraphOfAnabelioids.vertexFundamentalGroupImageEquiv}
  \lean{Iut.SourcePointedSemiGraphOfAnabelioids.edgeFundamentalGroupImageEquiv}
  \lean{Iut.SourcePointedSemiGraphOfAnabelioids.branchTransport}
  \lean{Iut.SourcePointedSemiGraphOfAnabelioids.subgroupDiagram}
  \lean{Iut.SourcePointedSemiGraphOfAnabelioids.subgroupDiagram_vertexGroup_isClosed}
  \lean{Iut.SourcePointedSemiGraphOfAnabelioids.subgroupDiagram_edgeGroup_isClosed}
  \lean{Iut.SourceSemiGraphSubgroupDiagram}
  \lean{Iut.SourceSemiGraphSubgroupDiagram.branchCosetMap}
  \lean{Iut.SourceSemiGraphSubgroupDiagram.branchCosetMap_smul}
  \lean{Iut.SourceSemiGraphSubgroupDiagram.cosetSemiGraph}
  \lean{Iut.SourceSemiGraphSubgroupDiagram.cosetAction}
  \lean{Iut.SourceSemiGraphSubgroupDiagram.stabilizer_leftCoset_eq_conjugate}
  \lean{Iut.SourceSemiGraphSubgroupDiagram.cosetAction_vertexStabilizer_eq_conjugate}
  \lean{Iut.SourceSemiGraphSubgroupDiagram.cosetAction_edgeStabilizer_eq_conjugate}
  \lean{Iut.SourceSemiGraphSubgroupDiagram.cosetAction_branchStabilizer_eq_edgeStabilizer}
  \lean{Iut.SourceProfiniteCosetSystem.levelSubgroup}
  \lean{Iut.SourceProfiniteCosetSystem.carrier}
  \lean{Iut.SourceProfiniteCosetSystem.system}
  \lean{Iut.SourceProfiniteCosetSystem.cosetSection}
  \lean{Iut.SourceProfiniteCosetSystem.iInf_levelSubgroup_eq}
  \lean{Iut.SourceProfiniteCosetSystem.sectionFiber}
  \lean{Iut.SourceProfiniteCosetSystem.sectionFiber_nonempty}
  \lean{Iut.SourceProfiniteCosetSystem.sectionFiber_isClosed}
  \lean{Iut.SourceProfiniteCosetSystem.sectionFiber_directed}
  \lean{Iut.SourceProfiniteCosetSystem.cosetSection_surjective}
  \lean{Iut.SourceProfiniteCosetSystem.cosetSection_injective}
  \lean{Iut.SourceProfiniteCosetSystem.cosetSectionEquiv}
  \lean{Iut.SourceProfiniteCosetSystem.sectionStabilizer_baseSection_eq}
  \lean{Iut.SourceProfiniteCosetSystem.sectionStabilizer_cosetSection_eq_conjugate}
  \lean{Iut.SourceSemiGraphSubgroupDiagram.profiniteVertexSection}
  \lean{Iut.SourceSemiGraphSubgroupDiagram.profiniteEdgeSection}
  \lean{Iut.SourceSemiGraphSubgroupDiagram.profiniteVertexSectionEquiv}
  \lean{Iut.SourceSemiGraphSubgroupDiagram.profiniteEdgeSectionEquiv}
  \lean{Iut.SourceSemiGraphSubgroupDiagram.profiniteVertexSection_stabilizer_eq}
  \lean{Iut.SourceSemiGraphSubgroupDiagram.profiniteEdgeSection_stabilizer_eq}
  \lean{Iut.SourceSemiGraphSubgroupDiagram.normalOpenLevel}
  \lean{Iut.SourceSemiGraphSubgroupDiagram.normalOpenLevelSemiGraph}
  \lean{Iut.SourceSemiGraphSubgroupDiagram.normalOpenLevel_vertices_finite}
  \lean{Iut.SourceSemiGraphSubgroupDiagram.normalOpenLevel_edges_finite}
  \lean{Iut.SourceSemiGraphSubgroupDiagram.normalOpenLevel_branchCosetMap_transition}
  \lean{Iut.SourceSemiGraphSubgroupDiagram.normalOpenLevelTransition}
  \lean{Iut.SourceSemiGraphSubgroupDiagram.normalOpenLevelVertexMap_comp}
  \lean{Iut.SourceSemiGraphSubgroupDiagram.normalOpenLevelEdgeMap_comp}
  \lean{Iut.SourceSemiGraphSubgroupDiagram.normalOpenLevelVertexMap_surjective}
  \lean{Iut.SourceSemiGraphSubgroupDiagram.normalOpenLevelEdgeMap_surjective}
  \lean{Iut.SourceSemiGraphSubgroupDiagram.normalOpenLevelVertexMap_action}
  \lean{Iut.SourceSemiGraphSubgroupDiagram.normalOpenLevelEdgeMap_action}
  \lean{Iut.SourceSemiGraphSubgroupDiagram.normalOpenLevelTransition_totalBranchMap_comp}
  \lean{Iut.SourceSemiGraphSubgroupDiagram.normalOpenLevelTransition_totalBranchMap_surjective}
  \lean{Iut.SourceSemiGraphSubgroupDiagram.normalOpenLevelTransition_totalBranchMap_action}
  \lean{Iut.SourceSemiGraphSubgroupDiagram.normalOpenLevelTransition_isProper}
  \lean{Iut.SourceGraphAction}
  \lean{Iut.SourceGraphAction.dist_image}
  \lean{Iut.SourceFiniteTree.induce_degree_ne_one_connected}
  \lean{Iut.SourceGraphAction.exists_invariant_subtree_card_le_two}
  \lean{Iut.SourceGraphAction.finite_fixesVertexOrEdge}
  \lean{Iut.SourceGraphAction.fixesVertexOrEdge}
  \lean{Iut.SourceSemiGraphTree}
  \lean{Iut.SourceSemiGraphTree.Action.fixedOpenEdge_of_fixed_boundary}
  \lean{Iut.SourceSemiGraphTree.Action.fixesVertexOrEdge}
  \lean{Iut.SourceGraphAction.fixes_path_pointwise}
  \lean{Iut.SourceGraphAction.subgroup_fixes_geodesic_pointwise}
  \lean{Iut.SourceSemiGraphTree.Action.FixesSubjoint}
  \lean{Iut.SourceSemiGraphTree.Action.fixesSubjoint_of_three_fixed_vertices}
  \lean{Iut.SourceSemiGraphTree.edge_abuts_vertex}
  \lean{Iut.SourceSemiGraphTree.Action.OverBase}
  \lean{Iut.SourceSemiGraphTree.Action.OverBase.not_swap_edge}
  \lean{Iut.SourceSemiGraphTree.Action.OverBase.fixedVertex_of_fixesVertexOrEdge}
  \lean{Iut.SourceCompactSemiGraphAction}
  \lean{Iut.SourceCompactSemiGraphAction.deckMap_range_finite}
  \lean{Iut.SourceCompactSemiGraphAction.fixesVertexOrEdge}
  \lean{Iut.SourceCompactSemiGraphAction.fixesVertex}
  \lean{Iut.SourceCompactSemiGraphAction.fixesSubjoint_of_three_fixed_vertices}
  \lean{Iut.SourceSemiGraphTree.Action.EquivariantLocalMap}
  \lean{Iut.SourceSemiGraphTree.Action.FixedSubjoint.map}
  \lean{Iut.SourceCofilteredFixedSubjointSystem}
  \lean{Iut.SourceCofilteredFixedSubjointSystem.exists_compatible_fixedSubjoints}
  \lean{Iut.SourceEstrangedIncidentBranchSystem}
  \lean{Iut.SourceEstrangedIncidentBranchSystem.subgroup_eq_bot_of_le_branch_intersection}
  \lean{Iut.SourceEstrangedIncidentBranchSystem.not_le_two_branches_of_ne_bot}
  \lean{Iut.SourceSemiGraphTree.Action.adjacent_of_fixedVertexSet_eq_pair}
  \lean{Iut.SourceSemiGraphTree.Action.FixedVertex.map}
  \lean{Iut.SourceCofilteredFixedVertexSystem}
  \lean{Iut.SourceCofilteredFixedVertexSystem.exists_compatible_fixedVertices}
  \lean{Iut.SourceCofilteredFixedVertexSystem.three_compatible_fixedVertices_not_pairwise_distinct}
  \lean{Iut.SourceIUTIProposition24PartI}
  \lean{Iut.SourceIUTIProposition24PartI.CuspidalInertia}
  \lean{Iut.SourceIUTIProposition24PartI.CuspidalInertia.conjugate_le_arithmeticTempered_iff}
  \lean{Iut.SourceIUTIProposition24PartI.CuspidalInertia.conjugate_classification}
  \lean{Iut.SourceIUTIProposition24PartI.CuspidalInertia.inertia_le_conjugate_arithmeticTempered_iff}
  \leanok
  The source semi-graph is first represented by a type of vertices, a type of
  edges, a two-element dependent branch fiber for each edge, and a partial
  coincidence map.  This retains loops, parallel edges, and closed, open, and
  isolated edges.  General morphisms carry branch-fiber equivalences and
  preserve only source abutments, exactly allowing an open branch to become
  verticial.  Lean constructs incident-branch maps, immersion and excision,
  proves that branchwise properness is equivalent to preservation of verticial
  cardinality, constructs sub-semi-graphs by truncating omitted abutments, and
  proves that compactification is a graph and constructs the compactification
  morphism induced by every general morphism.
  A coherent raw action is then defined on vertices, edges, and total branches.
  Fixing an incident branch forces fixation of its edge and abutting vertex,
  and every acting element permutes incident branches bijectively.  For a
  cofiltered finite system, Lean constructs the componentwise action on
  compatible sections and proves that the stabilizer of a compatible system is
  the intersection of all finite-level coordinate stabilizers.  This is the
  action-and-stabilizer portion of Remark~2.2.1; the graph-of-groups
  identification is then constructed on the Bass--Serre coset semi-graph.
  Branch transport elements allow general loop-edge and HNN-type incidence
  maps.  These maps are equivariant for the ambient action, and Lean proves
  that every lifted vertex and edge stabilizer is exactly the corresponding
  conjugate constituent subgroup.  For a closed constituent subgroup $H$ in a
  profinite ambient group $G$, open normal subgroups $N$ index actual finite
  coset levels $G/(H\vee N)$.  Lean constructs their equivariant transition
  functor and derives $\bigcap_N(H\vee N)=H$ from profinite separation.  For
  an arbitrary compatible system, its coordinate representative fibers are
  nonempty, closed, and directed; compactness produces one ambient element in
  their intersection.  Consequently $G/H$ is equivalent to the full type of
  compatible systems, and the compatible vertex and edge stabilizers equal
  the Bass--Serre stabilizers.
  For every profinite group $G$, Lean now constructs $\mathcal B(G)$ as the
  actual Galois category of finite discrete continuous $G$-actions and proves
  that its canonical fiber functor has fundamental group $G$.  Restriction
  along a continuous homomorphism is exact and induces that same homomorphism
  on certified fundamental groups.  Thus an open embedding $H\hookrightarrow
  G$ constructs the connected finite-etale open-image morphism
  $\mathcal B(H)\to\mathcal B(G)$.  Lean constructs the exact left-diagonal
  quotient convention of Proposition~1.2.1(v), proves it equivariantly
  equivalent to standard induction, and constructs the natural equivalence
  $\mathcal B(H)\simeq\mathcal B(G)_{G/H}$.  The source adjunction
  $j_{G/H}\dashv(-\times G/H)$ and an explicit induction--restriction
  adjunction identify restriction with product by $G/H$ under this
  equivalence.  For an arbitrary finite continuous $G$-set $S$, Lean
  constructs the conjugation action on dependent sections of every
  $T\to S$, proves the natural currying equivalence
  \[
    \operatorname{Hom}_{\mathcal B(G)_S}(S\times X,T)
      \simeq
    \operatorname{Hom}_{\mathcal B(G)}(X,\Gamma_S(T)),
  \]
  and obtains the second adjunction
  $(-\times S)\dashv\Gamma_S$.  Thus product by every $S$ preserves finite
  limits and finite colimits, proving Proposition~1.2.1(iii), and the
  pullback is packaged by the factorization criterion of
  Definition~1.2.2.  Lean also forms the finite orbit set $\pi_0(S)$,
  restricts an object over $S$ to every orbit, and glues an orbit-indexed
  family by dependent disjoint union.  Explicit natural isomorphisms in both
  directions prove
  \[
    \mathcal B(G)_S \simeq
      \prod_{c\in\pi_0(S)}\mathcal B(G)_{S_c}.
  \]
  For the canonical representative $s_c$, continuity makes
  $\operatorname{Stab}_G(s_c)$ open.  An equivariant orbit--stabilizer
  isomorphism and Proposition~1.2.1(v) then identify the $c$-factor with the
  certified connected anabelioid
  $\mathcal B(\operatorname{Stab}_G(s_c))$.  This proves
  Proposition~1.2.1(i) in the continuous-action model.  Lean also proves the
  categorical-equivalence prerequisite for the converse: from a pointed exact
  pullback equivalence it constructs the quasi-inverse pointed morphism, uses
  both triangle $2$-isomorphisms to prove bijectivity of the derived
  fundamental-group map up to the source's inner ambiguity, and obtains a
  continuous profinite-group equivalence.  It remains to prove that a
  connected slice object is transitive, construct the required comparison of
  basepoints, and classify the original morphism as restriction along the
  resulting open stabilizer.
  \lean{Iut.sourceDependentSection_continuous}
  \lean{Iut.sourceSliceDependentSectionHomEquiv}
  \lean{Iut.sourceSliceProductAdjDependentSection}
  \lean{Iut.sourceSliceProduct_preservesFiniteColimits}
  \lean{Iut.sourceSliceComponentProductEquivalence}
  \lean{Iut.sourceActionComponentCosetActionIso}
  \lean{Iut.sourceSliceAnabelioidEquivalence}
  \lean{Iut.sourceActionComponentEtaleFundamentalGroup}
  \lean{Iut.SourcePointedAnabelioidEquivalence.fundamentalGroupHom_bijective}
  \lean{Iut.SourcePointedAnabelioidEquivalence.fundamentalGroupContinuousMulEquiv}
  Each constituent anabelioid is now an actual Galois category with a chosen
  fiber functor and certified profinite fundamental group.  Branch and
  constituent-to-total morphisms are exact contravariant functors.  A
  categorical $2$-isomorphism derives its basepoint discrepancy, hence the
  conjugating branch transport, and Lean constructs the preceding closed
  subgroup diagram from the images of the induced fundamental-group maps.
  The constituent systems are assembled into normal-open coset semi-graphs.
  Normality extends every branch transport from $H_e$ to $H_e\vee N$; quotient
  maps give proper semi-graph morphisms with finite component fibers.  Lean
  proves incidence compatibility, surjectivity, identity and composition, and
  deck equivariance for their vertex, edge, and branch maps.  Constructing the
  total glued category $\mathcal B(\mathbb G)$ from the constituent categories
  and enriching the raw normal-open levels by connected finite etale
  component anabelioids remain open.

  A finite group acting on an arbitrary simple tree fixes a vertex or an edge
  setwise.  Lean proves this by constructing the finite connected invariant
  hull of one orbit and repeatedly pruning its invariant leaves until at most
  two vertices remain.  A group action by tree automorphisms also preserves
  graph distance; if a subgroup fixes two vertices, it fixes every vertex on
  their unique geodesic.  The source compactification is represented by marking
  the original vertices and requiring each appended boundary point to have a
  unique neighbor.  A fixed appended point therefore determines a fixed open
  edge.  Three distinct fixed original vertices also yield two fixed incident
  branches, hence a fixed subjoint.  This formalizes Lemma~1.8(ii)(a--c) in the
  marked-compactification encoding.  For the first step of Theorem~3.7(iii),
  a continuous compact-group action through the discrete deck group has finite
  image, so the preceding fixed-component theorem applies to the compact
  group.  Its oriented branch map over the base prevents an edge swap, while
  connectedness and the existence of an original vertex ensure that every
  edge abuts an original vertex.  Thus the compact group fixes an original
  vertex, as in the next sentence of the source proof.  At cofinal levels with
  three fixed vertices, the fixed-subjoint theorem no longer requires a finite
  acting group.  An equivariant locally injective quotient map transports the
  resulting subjoint to a finite level, and the finite cofiltered-limit theorem
  supplies a compatible system of such finite-level subjoints.  Definition~2.4(iv)
  is encoded on incident branch subgroups; Lean proves that containment in two
  estranged branch images forces the contained subgroup to be trivial.  Given
  the resulting cofinal two-point bounds, a second finite cofiltered-limit
  argument supplies a compatible fixed-vertex system, and a common-predecessor
  argument proves that there are at most two such systems.  If a level has
  exactly two fixed vertices, their unique geodesic is a single edge.  The exact subgroup
  configuration of \IUT{} I, Proposition~2.4(i), is typed separately.  From
  its compact-conjugate rigidity conclusion, normality of the geometric
  subgroup, and a nontrivial compact cusp inertia subgroup, Lean derives the
  two cusp-conjugacy conclusions of Corollary~2.5.  The associated topological
  realization and raw tree/action bridge, arbitrary finite semi-graph quotients, Remark~2.2.1's
  derivation of the two branch containments and the stabilizer-to-verticial
  subgroup identification, assembly of the cofinal pair bound from these
  pieces, maximal-compact/verticial classification, the
  inverse-limit proof of Proposition~2.4(i), and construction of cusp inertia
  from stable reduction remain open.
\end{definition}

\begin{definition}[Mono-theta environments and inverse-limit group]
  \label{def:mono-theta-environment-limit}
  \uses{def:initial-theta-data,def:tempered-semigraph-cusp-rigidity}
  \lean{Iut.TopologicalGroupCat}
  \lean{Iut.ContinuousAutomorphism.inner_range_normal}
  \lean{Iut.ContinuousOuterSubgroup}
  \lean{Iut.SourceMonoThetaEnvironment}
  \lean{Iut.SourceMonoThetaProjectiveSystem}
  \lean{Iut.SourceMonoThetaProjectiveSystem.inverseLimit}
  \lean{Iut.SourceIUTIIDefinition27PartI}
  \lean{Iut.SourceTopologicalSubquotient}
  \lean{Iut.SourceTopologicalSubquotient.projection}
  \lean{Iut.SourceTopologicalSubquotient.comap}
  \lean{Iut.SourceTopologicalSubquotient.comapToOriginal_injective}
  \lean{Iut.SourceTopologicalSubquotient.comapToOriginal_surjective_of_cuspidalInertia}
  \lean{Iut.SourceIUTIICorollary25CuspidalRestriction}
  \lean{Iut.SourceIUTIICorollary25CuspidalRestriction.restrictionMulEquiv}
  \lean{Iut.SourceIUTIICorollary25CuspidalRestriction.restrictionContinuousMulEquiv}
  \lean{Iut.SourceIUTIIDefinition27PartII}
  \lean{Iut.SourceIUTIIDefinition27PartII.exteriorCyclotomeSubquotient}
  \lean{Iut.SourceIUTIIDefinition27PartII.lDeltaThetaSubquotient}
  \lean{Iut.SourceIUTIIDefinition27PartII.restrictedTemperedSubquotient}
  \lean{Iut.SourceIUTIIDefinition27PartII.localGaloisSubquotient}
  \lean{Iut.SourceIUTIIDefinition27PartII.cyclotomicRigidity_source}
  \lean{Iut.SourceIUTIIDefinition27PartII.narrowInverseImage_le_wide}
  \lean{Iut.EtaleTheta.ThetaCenter}
  \lean{Iut.EtaleTheta.thetaCenter_le_center}
  \lean{Iut.EtaleTheta.thetaCenter_isMulCommutative}
  \lean{Iut.EtaleTheta.LDeltaTheta}
  \lean{Iut.EtaleTheta.lDeltaTheta_le_thetaCenter}
  \lean{Iut.EtaleTheta.lDeltaTheta_le_center}
  \lean{Iut.EtaleTheta.LDeltaThetaNumerator}
  \lean{Iut.EtaleTheta.LDeltaThetaDenominator}
  \lean{Iut.EtaleTheta.lDeltaThetaProjection_ker}
  \lean{Iut.EtaleTheta.lDeltaThetaProjection_surjective}
  \lean{Iut.EtaleTheta.lDeltaThetaSubquotientEquiv}
  \lean{Iut.GlobalLTorsionCoverInput.standardGeometricLDeltaThetaSubquotient}
  \lean{Iut.GlobalLTorsionCoverInput.standardGeometricLDeltaThetaMulEquiv}
  \lean{Iut.GlobalLTorsionCoverInput.standardArithmeticLDeltaThetaSubquotient}
  \lean{Iut.GlobalLTorsionCoverInput.standardGeometricToArithmeticLDeltaThetaMulEquiv}
  \lean{Iut.GlobalLTorsionCoverInput.standardArithmeticLDeltaThetaMulEquiv}
  \lean{Iut.GlobalLTorsionCoverInput.standardThetaRootArithmeticSubgroup_isOpen}
  \lean{Iut.GlobalLTorsionCoverInput.standardThetaRootSubquotient}
  \lean{Iut.GlobalLTorsionCoverInput.standardThetaRootContinuousMulEquiv}
  \lean{Iut.GlobalLTorsionCoverInput.standardModLThetaCenterSubquotient}
  \lean{Iut.GlobalLTorsionCoverInput.standardModLThetaCenterNumerator_eq_center_preimage}
  \lean{Iut.GlobalLTorsionCoverInput.standardModLThetaCenterContinuousMulEquivCenterImage}
  \lean{Iut.GlobalLTorsionCoverInput.standardThetaCenterImageEquiv}
  \lean{Iut.GlobalLTorsionCoverInput.standardModLThetaCenterMulEquiv}
  \lean{Iut.ProfiniteFundamentalExactSequence.localGaloisSubquotient}
  \lean{Iut.ProfiniteFundamentalExactSequence.localGalois_denominator_eq_geometric_range_comap}
  \lean{Iut.ProfiniteFundamentalExactSequence.localGaloisContinuousMulEquiv}
  \lean{Iut.ProfiniteFundamentalExactSequence.localGaloisContinuousMulEquiv_projection}
  \lean{Iut.SourceThetaFiniteLocalCoreData.xLocalGaloisSubquotient}
  \lean{Iut.SourceThetaFiniteLocalCoreData.xLocalGaloisContinuousMulEquiv}
  \leanok
  The source boundary for Etale Theta Definition~2.13 uses genuine
  topological groups, continuous automorphisms, the proved-normal exact range
  of inner automorphisms, the full inverse image of a subgroup of the outer
  automorphism group, a cyclic \(\mu_N\)-kernel, and
  the conjugacy class of the image of a continuous injective theta section.
  The IUT~II Definition~2.7 system is an actual functor with a certified
  categorical inverse limit; its ambient homomorphism has the exterior
  cyclotome as exact kernel and the tempered subgroup as exact image.  The two
  groups in Definition~2.7(ii) are computed as inverse images, and each of its
  four subquotients is represented by an actual topological quotient \(U/N\)
  by a normal subgroup.  Pullback is the derived quotient
  \(f^{-1}(U)/f^{-1}(N)\): the exterior cyclotome is restricted from the exact
  kernel in part (i), and the other three objects are pulled back from ambient
  subquotients.  Restricted cyclotomic rigidity is a continuous group
  equivalence between these derived endpoints.  For the local orbicurve exact
  sequence, the full arithmetic group modulo the geometric image is now
  constructed and proved topologically equivalent to the stored local Galois
  group.  Lean now distinguishes the integral and finite theta directions.
  The integral \(\Delta^\Theta\) is the kernel of the two-step-nilpotent
  quotient to the elliptic abelianization; a lower-central commutator
  calculation proves that this kernel is central and hence commutative.  The
  subgroup \(\ell\Delta^\Theta\) is the literal range of the \(\ell\)-th-power
  homomorphism on that center, and is consequently central and normal.  Pullback
  through \(\Delta_X\to\Delta_X^\Theta\) gives a literal numerator and
  lower-central denominator, and the first isomorphism theorem identifies the
  resulting geometric subquotient with \(\ell\Delta^\Theta\).  Its numerator
  and denominator are transported through the injective geometric-to-arithmetic
  profinite map, giving an actual arithmetic-image subquotient and a derived
  equivalence with the geometric coefficient group.  For a subgraph inclusion,
  the canonical map from the pulled-back quotient is proved injective; the two
  cusp-inertia facts cited in IUT~II, Corollary~2.5(i), derive its surjectivity
  and hence its algebraic and topological isomorphisms.  The isomorphism itself
  is not stored as an input.  Separately, the
  profinite theta-root cover is the proved-open inverse image of the selected
  Galois section in \(\Pi_X^{\mathrm{ell}}\).  The quotient
  \[
    \ker(\Pi_X\to\Pi_X^{\mathrm{ell}})/
    \ker(\Pi_X\to\Pi_X^\Theta)
  \]
  is recorded under its correct name: the finite mod-\(\ell\) theta-center
  residue.  It is identified topologically with the arithmetic exponent-\(\ell\)
  theta-center kernel and is not conflated with the integral subgroup
  \(\ell\Delta^\Theta\).  The abstract cusp-conjugacy consequences of
  \IUT{} I, Corollary~2.5 are now derived from Proposition~2.4(i).  Construction
  of the arithmetic tempered group and subgraph inclusion, derivation of
  Proposition~2.4(i) and the remaining generation/surjectivity facts from the
  stable-reduction semi-graph, identification of all three generic
  Definition~2.7 inputs, the full reduction morphisms, and the link to
  Corollary~2.8 remain source obligations.
\end{definition}

\begin{definition}[Continuous theta-value germs and Gaussian monoids]
  \label{def:continuous-theta-evaluation}
  \uses{def:mono-theta-environment-limit}
  \lean{Iut.ContinuousOneCocycle}
  \lean{Iut.ContinuousH1}
  \lean{Iut.ContinuousH1.toGerm}
  \lean{Iut.ContinuousH1Germ}
  \lean{Iut.ContinuousH1GermRepresentative.pullback}
  \lean{Iut.ContinuousH1Germ.restrict}
  \lean{Iut.ContinuousH1Germ.restrictMonoidHom}
  \lean{Iut.EquivariantOrbitMap}
  \lean{Iut.EquivariantOrbitMap.onOrbitQuotients}
  \lean{Iut.independentConjugacyMulAction}
  \lean{Iut.SourceIUTIIThetaEvaluationOrbits}
  \lean{Iut.SourceIUTIIIndependentConjugacyEvaluationAlgorithm}
  \lean{Iut.SourceIUTIIIndependentConjugacyEvaluationAlgorithm.toEquivariantOrbitMap}
  \lean{Iut.SourceIUTIIIndependentConjugacyEvaluationAlgorithm.onIndependentConjugacyClasses}
  \lean{Iut.SourceIUTIICorollary28RestrictionStages}
  \lean{Iut.SourceIUTIICorollary28RestrictionStages.afterSecond}
  \lean{Iut.SourceIUTIICorollary28RestrictionStages.totalRestrictionMonoidHom}
  \lean{Iut.SourceIUTIICorollary28RestrictionStages.afterSecond_zero_representatives}
  \lean{Iut.sourceConstantTorsion}
  \lean{Iut.sourceSplittingUpToTorsionMap}
  \lean{Iut.SourceSplittingUpToTorsion}
  \lean{Iut.sourceZeroLabelSplittingUpToTorsion}
  \lean{Iut.SourceIUTIIZeroLabelConstantComparison}
  \lean{Iut.SourceIUTIIZeroLabelConstantComparison.zeroToDiagonalMonoidHom}
  \lean{Iut.SourceIUTIIZeroLabelConstantComparison.zeroDiagonalSubmonoid_eq_diagonalRangeSubmonoid}
  \lean{Iut.SourceIUTIIZeroLabelConstantComparison.zeroToDiagonalMulEquiv}
  \lean{Iut.SourceIUTIIZeroLabelConstantComparison.zeroToDiagonalContinuousMulEquiv}
  \lean{Iut.SourceIUTIIZeroLabelConstantComparison.zeroToDiagonalMonoidHom_equivariant}
  \lean{Iut.etaleThetaUnitActionOnSignLabel}
  \lean{Iut.iutIISignQuotientPermHom}
  \lean{Iut.iutIISignQuotientPermHom_transitive}
  \lean{Iut.iutIISignQuotientAction_transitive}
  \lean{Iut.SourceTopologicalGroupMonoidActionPair}
  \lean{Iut.SourceTopologicalGroupMonoidActionPair.Iso}
  \lean{Iut.SourceTopologicalGroupMonoidActionPair.Iso.castTarget}
  \lean{Iut.SourceIUTIISymmetrizingIsomorphisms}
  \lean{Iut.SourceIUTIISymmetrizingIsomorphisms.outerTransport_one}
  \lean{Iut.SourceIUTIISymmetrizingIsomorphisms.outerTransport_mul}
  \lean{Iut.SourceIUTIISymmetrizingIsomorphisms.fromBase}
  \lean{Iut.SourceIUTIISymmetrizingIsomorphisms.diagonalMonoidHom}
  \lean{Iut.SourceIUTIISymmetrizingIsomorphisms.diagonalMonoidHom_injective}
  \lean{Iut.SourceIUTIISymmetrizingIsomorphisms.diagonalRangeSubmonoid}
  \lean{Iut.SourceIUTIISymmetrizingIsomorphisms.diagonalRangeContinuousMulEquiv}
  \lean{Iut.SourceIUTIISymmetrizingIsomorphisms.diagonalRangeContinuousMulEquiv_equivariant}
  \lean{Iut.SourceIUTIISymmetrizingIsomorphisms.diagonalUnitSubmonoid}
  \lean{Iut.SourceIUTIISymmetrizingIsomorphisms.diagonalActionHom}
  \lean{Iut.SourceIUTIISymmetrizingIsomorphisms.diagonalMonoidHom_equivariant}
  \lean{Iut.SourceIUTIISymmetrizingIsomorphisms.gaussianMonoid}
  \lean{Iut.SourceIUTIISymmetrizingIsomorphisms.profile_mem_gaussianMonoid}
  \lean{Iut.SourceIUTIISymmetrizingIsomorphisms.gaussianMonoidActionHom}
  \lean{Iut.MulEquivsDifferBySingleInner}
  \lean{Iut.singleInnerIsomorphismSetoid}
  \lean{Iut.SourceSingleInnerIsomorphismClass}
  \lean{Iut.SourceIUTIISynchronizedRestrictionIsomorphismClass}
  \lean{Iut.SourceIUTIICorollary35RestrictionContext}
  \lean{Iut.SourceIUTIICorollary35RestrictionDiagram}
  \lean{Iut.SourceIUTIICorollary35SingleBasepointAction}
  \lean{Iut.RestrictionDiagramsDifferBySingleInner}
  \lean{Iut.singleInnerRestrictionDiagramSetoid}
  \lean{Iut.SourceIUTIISynchronizedRestrictionDiagramClass}
  \lean{Iut.SourceIUTIICorollary35RestrictionDiagram.synchronizedClassToMonoidIsomorphismClass}
  \lean{Iut.SourceIUTIIValueProfile}
  \lean{Iut.TorsionUnitSubgroup}
  \lean{Iut.SourceIUTIIMultiplicativeTorsionAction}
  \lean{Iut.SourceIUTIIMultiplicativeTorsionAction.actionMap_injective}
  \lean{Iut.SourceIUTIIMultiplicativeTorsionAction.toTorsionUnitSubgroup}
  \lean{Iut.SourceIUTIIMultiplicativeTorsionAction.toTorsionThetaValueFamily}
  \lean{Iut.SourceIUTIIMultiplicativeTorsionAction.thetaValues_eq_orbitCarriers}
  \lean{Iut.SourceIUTIIMultiplicativeTorsionAction.toFiniteOrderAction}
  \lean{Iut.SourceIUTIIMultiplicativeFiniteOrderAction}
  \lean{Iut.SourceIUTIIMultiplicativeFiniteOrderAction.orbitCarrier_subset_torsion}
  \lean{Iut.SourceIUTIITorsionThetaValueFamily}
  \lean{Iut.SourceIUTIIInfiniteThetaRelation}
  \lean{Iut.SourceTorsionCyclotomicCommMonoid}
  \lean{Iut.SourceTorsionCyclotomicTopologicalCommMonoid}
  \lean{Iut.SourceMonoidCyclotome}
  \lean{Iut.SourceMLFAlgebraicIntegerRing}
  \lean{Iut.SourceMLFIntegralMonoid}
  \lean{Iut.SourceMLFAbsoluteGaloisGroup}
  \lean{Iut.SourceMLFIntegralMonoid.galoisAction}
  \lean{Iut.SourceMLFIntegralMonoid.algebraicClosure_isFractionRing}
  \lean{Iut.SourceMLFIntegralMonoid.groupificationEquivAlgebraicClosureUnits}
  \lean{Iut.SourceMLFIntegralMonoid.torsionUnit}
  \lean{Iut.SourceMLFIntegralMonoid.torsionUnit_image}
  \lean{Iut.SourceMLFCyclotome.rationalCircleToComplexTorsion_bijective}
  \lean{Iut.SourceMLFCyclotome.torsionUnitsMap_bijective}
  \lean{Iut.SourceMLFCyclotome.rationalCircleMulEquivAlgebraicClosureUnitsTorsion}
  \lean{Iut.SourceMLFIntegralMonoid.torsionToAlgebraicClosureTorsion_bijective}
  \lean{Iut.SourceMLFIntegralMonoid.torsionEquivAlgebraicClosureTorsion}
  \lean{Iut.SourceMLFIntegralMonoid.groupificationRootableByNat}
  \lean{Iut.SourceTorsionCyclotomicCommMonoid.rootsOfUnity_finite}
  \lean{Iut.SourceTorsionCyclotomicCommMonoid.rationalCircle_torsion_natCard}
  \lean{Iut.SourceTorsionCyclotomicCommMonoid.rootsOfUnityEquiv}
  \lean{Iut.SourceTorsionCyclotomicCommMonoid.rootsOfUnity_natCard}
  \lean{Iut.SourceTorsionCyclotomicCommMonoid.RootsOfUnityGroup}
  \lean{Iut.SourceTorsionCyclotomicCommMonoid.rootsOfUnityGroup_natCard}
  \lean{Iut.SourceKummerFaithfulness.eq_one_of_roots_of_residuallyFinite}
  \lean{Iut.SourceKummerFaithfulness.residuallyFinite_of_profinite}
  \lean{Iut.SourceKummerFaithfulness.eq_one_of_roots_in_profinite}
  \lean{Iut.SourceMLFKummerFaithfulness.finite_quotient_map}
  \lean{Iut.SourceMLFKummerFaithfulness.units_residuallyFinite_of_finite_quotients}
  \lean{Iut.SourceMLFKummerFaithfulness.integralClosureUnits_residuallyFinite}
  \lean{Iut.SourceMonoidCyclotome.topology}
  \lean{Iut.SourceMonoidCyclotome.boundedCoe_isClosedEmbedding}
  \lean{Iut.SourceMonoidCyclotome.asProfinite}
  \lean{Iut.SourceMonoidCyclotome.eq_one_of_roots}
  \lean{Iut.SourceMLFModelTMStructure}
  \lean{Iut.SourceMLFModelTMStructure.canonical}
  \lean{Iut.SourceMLFModelTMPair}
  \lean{Iut.SourceMLFModelTMPair.monoAnalytic}
  \lean{Iut.SourceMLFModelTMPair.monoAnalytic_augmentation_apply}
  \lean{Iut.SourceMLFGaloisTMPairData}
  \lean{Iut.SourceMLFGaloisTMPairData.cyclotomeAction}
  \lean{Iut.SourceMLFGaloisTMPair}
  \lean{Iut.SourceMLFGaloisTMPair.monoAnalytic}
  \lean{Iut.SourceMLFGaloisTMPair.monoAnalytic_augmentation_apply}
  \lean{Iut.SourceMLFGaloisComparison.continuous_autCongr}
  \lean{Iut.SourceMLFGaloisComparison.continuousAutCongr}
  \lean{Iut.SourceMLFGaloisComparison.separableClosureEquivAlgebraicClosure}
  \lean{Iut.SourceMLFGaloisComparison.absoluteGaloisEquiv}
  \lean{Iut.SourceMLFGaloisTMPair.monoAnalyticSeparable}
  \lean{Iut.SourceMLFGaloisTMPair.pullbackActingGroup}
  \lean{Iut.SourceThetaFiniteLocalCoreData.mlfGaloisTMPair}
  \lean{Iut.SourceThetaFiniteLocalCoreData.unitKummerEmbedding}
  \lean{Iut.SourceThetaFiniteLocalCoreData.unitKummer_injective}
  \lean{Iut.SourceThetaFiniteLocalCoreData.xLocalGaloisTMPair}
  \lean{Iut.SourceThetaFiniteLocalCoreData.xLocalGaloisTMPair_action}
  \lean{Iut.SourceThetaFiniteLocalCoreData.xLocalGaloisCoefficientAction}
  \lean{Iut.SourceThetaFiniteLocalCoreData.xLocalGaloisCoefficientAction_act}
  \lean{Iut.SourceThetaFiniteLocalCoreData.xLocalGaloisH1Restriction}
  \lean{Iut.SourceThetaFiniteLocalCoreData.xLocalGaloisRestrictedMonoidKummer}
  \lean{Iut.SourceThetaFiniteLocalCoreData.xLocalGaloisRestrictedMonoidKummer_commutes}
  \lean{Iut.SourceThetaFiniteLocalCoreData.xLocalGaloisRestrictedUnitKummer}
  \lean{Iut.SourceThetaFiniteLocalCoreData.xLocalGaloisRestrictedUnitKummer_commutes}
  \lean{Iut.SourceThetaFiniteLocalCoreData.xLocalGaloisUnitKummerEmbedding}
  \lean{Iut.SourceThetaFiniteLocalCoreData.xLocalGaloisUnitKummer_injective}
  \lean{Iut.SourceFiniteTimesMuKummerInput}
  \lean{Iut.SourceFiniteTimesMuKummerInput.timesMuMonoid}
  \lean{Iut.SourceFiniteTimesMuKummerInput.timesMuUnitsEquiv}
  \lean{Iut.SourceFiniteTimesMuKummerFrobenioid}
  \lean{Iut.SourceFiniteTimesMuKummerFrobenioid.rationalUnitsEquiv}
  \lean{Iut.SourceArchimedeanKummerSystem.unitQuotientPresentation}
  \lean{Iut.SourceArchimedeanKummerSystem.noncompact_eq_one_of_isUnit}
  \lean{Iut.SourceArchimedeanKummerSystem.quotientSplitTopologicalMonoid}
  \lean{Iut.SourceArchimedeanKummerSystem.quotientSplitTransition}
  \lean{Iut.SourceArchimedeanKummerSystem.FrobenioidStage}
  \lean{Iut.SourceArchimedeanKummerSystem.FrobenioidSystem}
  \lean{Iut.SourceArchimedeanKummerSystem.FrobenioidSystem.divisor_compatible}
  \lean{Iut.SourceArchimedeanKummerSystem.FrobenioidSystem.frobeniusDegree_compatible}
  \lean{Iut.SourceArchimedeanKummerSystem.FrobenioidSystem.carrierDiagram}
  \lean{Iut.SourceArchimedeanKummerSystem.FrobenioidSystem.baseDiagram}
  \lean{Iut.SourceArchimedeanKummerSystem.FrobenioidSystem.rationalMonoidDiagram}
  \lean{Iut.SourceArchimedeanKummerSystem.MultiplicativelyCofinalSubset.indexInclusionFinal}
  \lean{Iut.SourceArchimedeanKummerSystem.MultiplicativelyCofinalSubset.isColimitEquivRestricted}
  \lean{Iut.SourceMLFGaloisTMPair.grothendieckAction}
  \lean{Iut.SourceMLFGaloisTMPair.fixedSubmonoid}
  \lean{Iut.SourceMLFIntegralMonoid.absoluteOpenStabilizer}
  \lean{Iut.SourceMLFGaloisTMPair.openStabilizer}
  \lean{Iut.SourceMLFGaloisTMPair.openStabilizer_fixed}
  \lean{Iut.SourceMLFGaloisTMPair.continuousCyclotomeAction}
  \lean{Iut.SourceMLFGaloisTMPair.rootableRationalPowers}
  \lean{Iut.SourceMLFGaloisTMPair.groupificationRootableByNat}
  \lean{Iut.SourceMLFGaloisTMPair.CompatibleRootSystem}
  \lean{Iut.SourceMLFGaloisTMPair.CompatibleRootSystem.ofModel}
  \lean{Iut.SourceMLFGaloisTMPair.CompatibleRootSystem.quotient_pow_den}
  \lean{Iut.SourceMLFGaloisTMPair.CompatibleRootSystem.galoisTranslate}
  \lean{Iut.SourceMLFGaloisTMPair.CompatibleRootSystem.quotientUnit_exists}
  \lean{Iut.SourceMLFGaloisTMPair.KummerRootRealization}
  \lean{Iut.SourceMLFGaloisTMPair.KummerRootRealization.ofRootSystem}
  \lean{Iut.SourceMLFGaloisTMPair.KummerRootRealization.ratioCyclotome}
  \lean{Iut.SourceMLFGaloisTMPair.KummerRootRealization.ratioCyclotome_continuous}
  \lean{Iut.SourceMLFGaloisTMPair.ContinuousKummerRootRealization}
  \lean{Iut.SourceMLFGaloisTMPair.ContinuousKummerRootRealization.cocycle}
  \lean{Iut.SourceMLFGaloisTMPair.CompatibleRootComparison}
  \lean{Iut.SourceMLFGaloisTMPair.CompatibleRootComparison.comparisonUnit_exists}
  \lean{Iut.SourceMLFGaloisTMPair.CompatibleRootComparison.ofModel}
  \lean{Iut.SourceMLFGaloisTMPair.CompatibleRootComparison.germ_eq}
  \lean{Iut.SourceMLFGaloisTMPair.CompatibleRootComparison.localClass_eq}
  \lean{Iut.SourceMLFGaloisTMPair.ContinuousModelKummerRootChoice}
  \lean{Iut.SourceMLFGaloisTMPair.KummerRootTheory}
  \lean{Iut.SourceMLFGaloisTMPair.KummerRootTheory.ofContinuousEquiv}
  \lean{Iut.SourceMLFGaloisTMPair.KummerRootTheory.chosen_rootSystem}
  \lean{Iut.SourceMLFGaloisTMPair.KummerRootTheory.kummerMap}
  \lean{Iut.SourceMLFGaloisTMPair.KummerRootTheory.canonicalKummerMap}
  \lean{Iut.SourceMLFGaloisTMPair.ContinuousLocalModelKummerRootRealization}
  \lean{Iut.SourceMLFGaloisTMPair.LocalKummerRootTheory}
  \lean{Iut.SourceMLFGaloisTMPair.LocalKummerRootTheory.ofContinuousEquiv}
  \lean{Iut.SourceMLFGaloisTMPair.LocalKummerRootTheory.chosen_rootSystem}
  \lean{Iut.SourceMLFGaloisTMPair.LocalKummerRootTheory.kummerMap}
  \lean{Iut.SourceMLFGaloisTMPair.LocalKummerRootTheory.local_toGerm_eq_global}
  \lean{Iut.SourceMLFGaloisTMPair.LocalKummerRootTheory.canonicalKummerMap}
  \lean{Iut.SourceMLFGaloisTMPair.LocalKummerRootTheory.canonicalKummerMap_toGerm}
  \lean{Iut.SourceMLFGaloisTMPair.augmentation}
  \lean{Iut.SourceMLFGaloisTMPair.augmentation_surjective}
  \lean{Iut.SourceMLFGaloisTMPair.augmentation_image_isOpen}
  \lean{Iut.SourceMLFGaloisTMPair.augmentationOpenImage}
  \lean{Iut.SourceMLFGaloisTMPair.fixedFieldOfOpenSubgroup}
  \lean{Iut.SourceMLFGaloisTMPair.fixedFieldOfOpenSubgroup_finiteDimensional}
  \lean{Iut.SourceMLFGaloisTMPair.modelGroupificationEquiv}
  \lean{Iut.SourceMLFGaloisTMPair.modelGroupificationEquiv_action}
  \lean{Iut.SourceMLFGaloisTMPair.fixedFieldIntegralUnit}
  \lean{Iut.SourceIUTIIUnitKummerEmbedding}
  \lean{Iut.SourceIUTIIUnitKummerEmbedding.canonical}
  \lean{Iut.SourceIUTIIUnitKummerEmbedding.monoAnalytic}
  \lean{Iut.SourceIUTIIUnitKummerEmbedding.canonical_monoidKummer}
  \lean{Iut.SourceIUTIIUnitKummerEmbedding.canonical_unitKummer}
  \lean{Iut.SourceIUTIIUnitKummerEmbedding.monoidKummer}
  \lean{Iut.SourceIUTIIUnitKummerEmbedding.unitKummer}
  \lean{Iut.SourceIUTIIUnitKummerEmbedding.fixed_roots_of_unitKummer_eq_one}
  \lean{Iut.SourceIUTIIUnitKummerEmbedding.unit_eq_one_of_unitKummer_eq_one}
  \lean{Iut.SourceIUTIIUnitKummerEmbedding.unitKummer_injective}
  \lean{Iut.SourceIUTIIUnitKummerEmbedding.canonical_unitKummer_injective}
  \lean{Iut.SourceIUTIIUnitKummerEmbedding.monoAnalytic_unitKummer_injective}
  \lean{Iut.SourceIUTIIUnitKummerEmbedding.constantSubgroup}
  \lean{Iut.SourceIUTIIUnitKummerEmbedding.torsionUnitImage}
  \lean{Iut.SourceIUTIIUnitKummerEmbedding.torsionUnitImage_eq_constant_torsion}
  \lean{Iut.SourceIUTIIThetaEvaluationOrbits.toTorsionThetaValueFamily}
  \lean{Iut.SourceIUTIIThetaEvaluationOrbits.toTorsionThetaValueFamily_thetaValues}
  \lean{Iut.SourceIUTIITorsionThetaValueFamily.representativeProfile}
  \lean{Iut.SourceIUTIITorsionThetaValueFamily.profiles_differByPointwiseTorsion}
  \lean{Iut.SourceIUTIIValueProfile.equivPi}
  \lean{Iut.SourceIUTIIValueProfile.natCard_eq_twoEll_pow_lStar}
  \lean{Iut.sourceGaussianMonoid}
  \lean{Iut.ValueProfilesDifferByPointwiseTorsion}
  \lean{Iut.sourceTwoEllGaussianSubmonoid}
  \lean{Iut.sourceTwoEllGaussianSubmonoid_eq}
  \lean{Iut.SourceIUTIITorsionThetaValueFamily.twoEllGaussianSubmonoid_eq}
  \lean{Iut.SourceIUTIITorsionThetaValueFamily.profileIndependentTwoEllGaussianSubmonoid}
  \lean{Iut.SourceIUTIITorsionThetaValueFamily.twoEllGaussianSubmonoid_eq_profileIndependent}
  \lean{Iut.ValueProfileRootSystemsDifferByDiagonalTorsion}
  \lean{Iut.sourceInfiniteGaussianMonoid}
  \lean{Iut.sourceInfiniteGaussianMonoid_eq_of_differByDiagonalTorsion}
  \lean{Iut.sourceConstantTorsion_unitKummer}
  \lean{Iut.sourceConstantTorsion_map}
  \lean{Iut.sourceConstantTorsionQuotientCongr}
  \lean{Iut.SourceSplittingUpToTorsion.map}
  \lean{Iut.SourceFiniteGaussianSplittingUpToTorsion}
  \lean{Iut.SourceInfiniteGaussianSplittingUpToTorsion}
  \lean{Iut.sourceGaussianMonoid_carrier_eq_splittingFactorProduct}
  \lean{Iut.sourceInfiniteGaussianMonoid_carrier_eq_splittingFactorProduct}
  \lean{Iut.sourceCorollary35FiniteGaussianSplittingOfProposition31}
  \lean{Iut.sourceCorollary35InfiniteGaussianSplittingOfProposition31}
  \lean{Iut.SourceTorsionDetectingRestriction}
  \lean{Iut.SourceTorsionDetectingRestriction.splittingUpToTorsion}
  \lean{Iut.SourceIUTIIPointedInversion}
  \lean{Iut.SourceIUTIIPointedInversion.decompositionInclusion}
  \lean{Iut.SourceIUTIICorollary112PointedRestriction}
  \lean{Iut.SourceIUTIICorollary112PointedRestriction.detects_constant_torsion}
  \lean{Iut.SourceIUTIICorollary112PointedRestriction.toTorsionDetectingRestriction}
  \lean{Iut.SourceIUTIICorollary112PointedRestriction.splittingUpToTorsion}
  \lean{Iut.sourceCorollary35FiniteGaussianSplittingOfRestrictionData}
  \lean{Iut.sourceCorollary35InfiniteGaussianSplittingOfRestrictionData}
  \lean{Iut.SourceIUTIICorollary28RestrictionStages.totalRestriction_mem_afterSecond_theta_carrier}
  \lean{Iut.SourceIUTIICorollary28RestrictionStages.totalRestriction_mem_afterSecond_infiniteTheta_carrier}
  \lean{Iut.SourceIUTIICorollary28RestrictionStages.zeroThetaTorsionDetectingRestriction}
  \lean{Iut.SourceIUTIICorollary28RestrictionStages.zeroInfiniteThetaTorsionDetectingRestriction}
  \lean{Iut.SourceIUTIICorollary28RestrictionStages.zeroThetaSplittingUpToTorsion}
  \lean{Iut.SourceIUTIICorollary28RestrictionStages.zeroInfiniteThetaSplittingUpToTorsion}
  \lean{Iut.thetaPowerValueProfile}
  \lean{Iut.thetaPowerGaussianMonoid}
  \lean{Iut.SourceIUTIIHodgeArakelovThetaPilot}
  \lean{Iut.SourceIUTIIHodgeArakelovThetaPilot.thetaPowerProfile}
  \lean{Iut.SourceIUTIIHodgeArakelovThetaPilot.finiteGaussianMonoid}
  \lean{Iut.SourceIUTIIHodgeArakelovThetaPilot.infiniteGaussianMonoid}
  \lean{Iut.SourceIUTIIHodgeArakelovThetaPilot.twoEllGaussianSubmonoid_eq_profileIndependent}
  \lean{Iut.SourceIUTIIHodgeArakelovThetaPilot.thetaOrbit_natCard}
  \lean{Iut.SourceIUTIIHodgeArakelovThetaPilot.valueProfiles_natCard}
  \lean{Iut.SourceIUTIIHodgeArakelovThetaPilot.zeroTheta_carrier_subset_torsion}
  \leanok
  Continuous nonabelian \(H^1\) is defined as continuous crossed
  homomorphisms modulo coboundaries.  The displayed direct limit over open
  subgroups is an explicit quotient of cocycle germs.  Pullback along a
  continuous group homomorphism is constructed on representatives and proved
  to descend to the quotient.  For abelian coefficients, pointwise
  multiplication descends to commutative group structures on both continuous
  \(H^1\) and its germ colimit, and restriction is a homomorphism.  Given the
  two paper homomorphisms and
  equivariance with the \(\mu_{2\ell}\)- and \(\mu\)-actions, the first and
  second restricted theta-orbit families are derived rather than stored as
  output fields.  Identity representatives at the zero label survive both
  restrictions, and every ambient orbit value is proved to map into the
  corresponding final orbit.  The inertia choice and ambient-subgroup choice carry
  independent copies of the conjugacy action; a two-variable equivariant
  evaluation descends to the corresponding product-action orbit quotients.
  The defining infinite-theta condition is stated exactly: membership means
  that some positive power agrees with a theta value up to torsion.  At the
  constant boundary, the model arithmetic monoid is literally the nonzero
  part of the integral closure of the MLF valuation ring in its algebraic
  closure.  Its Krull-topological absolute Galois group and algebraic action
  are constructed, a model pair has a continuous surjective Galois
  augmentation, and every general pair is required to be isomorphic to such
  a model.  Following Absolute Anabelian Topics~III, Remark~3.1.1, the exact
  \(T_{\rm M}\) certificate uses the equivalent topology-free monoid
  presentation; an arbitrary topology therefore cannot certify a model.
  The torsion-cyclotomic structure is no longer supplied as a certificate.
  Lean identifies the rational circle with all finite-order complex units,
  proves that embeddings between algebraically closed characteristic-zero
  fields are bijective on \(n\)-th roots by their exact cardinalities, and
  proves that every finite-order element of
  \(\overline{k}^{\times}\) is an integral unit.  Thus the identification
  \(O_{\overline{k}}^{\triangle,\mu}\cong\mathbb Q/\mathbb Z\), the
  identity-augmentation mono-analytic model pair, and its coefficient-identity
  unit-Kummer embedding are source-constructed without caller data.
  The fraction-field theorem identifies the actual Grothendieck group of this
  cancellative monoid with \(\overline{k}^{\times}\).  Algebraic closedness
  makes this group rootable, and injectivity of divisible abelian groups
  extends the prescribed map \(\mathbb Z\to\overline{k}^{\times}\) across
  \(\mathbb Z\to\mathbb Q\).  Thus coherent compatible rational roots are
  now constructed from the literal model and transported to every model pair;
  Kummer theories must realize those roots rather than supply arbitrary ones.
  Two systems agree on integer exponents, so their quotient at \(q\) has order
  dividing the denominator of \(q\).  This finite-order element and its inverse
  are integral, hence construct an actual unit of
  \(O_{\overline{k}}^{\triangle}\); transport through the model isomorphism
  constructs every root-system comparison rather than assuming it.
  Applying Galois to the constructed root system and comparing it with the
  original constructs every integral Galois root ratio from a fixed-point
  proof.  Krull stabilizers of the underlying algebraic integers are open and
  transport to the general pair.  The root ratios are locally constant in
  discrete units, hence continuous in the pointwise cyclotome topology.  The
  latter is a closed subspace of a product of finite root-of-unity groups, so
  Lean proves that it is compact, Hausdorff, totally disconnected, and bundles
  it as a profinite space.  The ratios are proved multiplicative, killed
  on \(\mathbb Z\), descended through \(\mathbb Q/\mathbb Z\), and assembled
  into continuous cocycles.  Comparing two root systems gives the explicit
  coboundary between their cocycles.  Hence both Kummer maps
  \(M^H\to H^1(H,\mu_{\mathbb Z}(M))\) and
  \(M\to\varinjlim_JH^1(J,\mu_{\mathbb Z}(M))\) are derived monoid
  homomorphisms, and the canonical local map is proved compatible with the
  canonical germ map.
  The unit Kummer homomorphism is its restriction, never an independently
  supplied field; injectivity proves that \(M_{\rm TM}^{\mu}\) is exactly the
  image of the torsion units.  Proposition~3.2(ii) constructs the map but does
  not assert injectivity.  Following Definition~1.5 and Remark~1.5.4(i), Lean
  now proves that the integral-closure unit group in every finite MLF extension
  is residually finite: Krull intersection separates a nontrivial unit, while
  reduction modulo a suitable power of the base maximal ideal gives a finite
  quotient.  For a compact IUT acting group, Lean also proves that the image of
  each open germ domain under the arithmetic augmentation is open: it has
  finite index, and compactness makes its image closed.  Infinite Galois
  correspondence then constructs the associated finite fixed extension.
  Equality of a constructed unit Kummer germ with one is now unpacked into a
  common open subgroup and an explicit coboundary.  Evaluating its cochain at
  (1/n) adjusts the compatible root; Lean proves that the adjusted root has
  (n)-th power equal to the original unit and is fixed by the common
  subgroup.  Transport to the finite fixed extension makes the root and its
  inverse integral by the power criterion.  Residual finiteness of that
  integral-unit group forces the original unit to be one, so the constructed
  unit Kummer homomorphism is injective without a stored faithfulness field.
  The local endpoint is now the exact one used by Initial Theta data.  Lean
  proves continuity of Galois conjugation across an algebra equivalence by
  pulling each finite-intermediate-field stabilizer back to the stabilizer of
  its inverse image.  In characteristic zero, the separable closure therefore
  gives a topological absolute-Galois group equivalent to the algebraic-closure
  model.  The canonical valuation on every finite number-field completion is
  proved compatible with its topology and makes the completion an MLF.  These
  results transport the canonical pair and its injective coefficient-identity
  unit Kummer map to the literal local absolute Galois group stored by each
  finite Initial Theta core.  Exactness then identifies that group with the
  arithmetic fundamental-group quotient by the geometric image.  Lean
  transports the MLF pair and faithful canonical unit Kummer map to this
  literal quotient, pulls back the cyclotomic coefficient action, constructs
  the induced continuous-\(H^1\) restriction, and proves the monoid and unit
  Kummer restriction squares commute.  Constructing the restricted-tempered
  mono-theta map to this endpoint remains open.
  At a finite place, the Definition~4.9 input now ties the selected splitting
  at an isotropic reference object to an actual MLF-Galois \(TM\)-pair, a
  stable characteristic submonoid, and the genuine \(\operatorname{Ism}\)
  Kummer orbit.  Lean constructs the bad-place quotient by
  \(\mu_{2\ell}\) and the good-place characteristic factor, forms the literal
  product with \(O^{\times\mu}\), and proves that this is its complete unit
  group.  The reconstruction output must have the original base category and
  this exact reference rational monoid, so an unrelated Frobenioid cannot
  inhabit the interface.  Derivation of that output from Frobenioids~I,
  Theorem~5.2(ii), remains explicit.
  At an archimedean place, the exact \(N\)-stage replaces the compact factor by
  \(S^1/\mu_N\) and retains the original noncompact factor.  Lean proves that
  this noncompact factor has no nontrivial units and constructs the continuous
  transition which is the torsion-quotient map on the compact factor and the
  identity on the noncompact factor.  A reconstructed stage must be an actual
  split Frobenioid whose split topological monoid is structure-preservingly
  isomorphic to this exact product and whose base is equivalent to the original
  base.  Its transition functors preserve the base projection, FSM morphisms,
  divisors, and Frobenius degrees.  Carrier categories, bases, and reference
  rational monoids then form strict divisibility-indexed diagrams.  Every
  multiplicatively cofinal restriction is proved final, so it preserves and reflects the universal
  colimit cocone.  Construction of these stages and transition functors from
  the original archimedean \(TM\)/Frobenioid data remains explicit.
  Definition~4.9(vi)--(viii) is now represented by a universe-polymorphic
  \(F^{\vdash\times\mu}\)-prime-strip boundary.  It fixes the bad/good
  finite construction at every selected finite place and the complete
  torsion-quotient system at every archimedean place.  The corresponding
  \(F^\times\)- and \(F^{\times\mu}\)-components are the actual wide
  subcategories of divisor-zero isometries.  The global record retains the
  realified Frobenioid, prime/place bijection, and every \(\rho_v\).  Its pilot
  is computed by transporting typed cyclic bad-place generators through these
  \(\rho_v\)'s and applying a global reconstruction function, and its
  arithmetic degree is required to be negative.  Supplying this reconstruction
  from Corollaries~4.6 and~4.10 remains open.  At each finite place, the
  structure-preserving isomorphism layer now transports the topological Galois
  group, arithmetic action, unit-mod-torsion quotient, every open-subgroup
  invariant Kummer image, both product generators, the reconstructed split
  Frobenioid, its base, reference object, and rational monoid.  Identity and
  composition are constructed, and preservation of the genuine
  \(\operatorname{Ism}\)-orbit is derived.  At an archimedean place, the
  system isomorphism transports the original split Frobenioid, every torsion
  quotient, quotient and orientation maps, level transitions, reconstructed
  stages, bases, divisors, Frobenius degrees, FSM morphisms, split products,
  and rational transitions.  Identity and composition are constructed.  The
  level maps restrict to the isometry coarsifications and form a natural
  transformation of complete divisibility diagrams with equivalence
  components.  The assembled globally realified isomorphism now combines
  these place-indexed local equivalences with a
  structure-preserving equivalence of the global Frobenioids, the prime/place
  and \(\rho_v\) squares, bad-place generator and currency compatibility,
  reconstruction for every character tuple, and arithmetic-degree
  compatibility.  Identity and composition are constructed.  Compatibility
  of the distinguished pilot character, pilot object, and pilot degree is
  derived.  Following IUT~I, Section~0, the place-indexed and globally
  realified full poly-isomorphisms are literal quotients of all complete
  representatives.  Categorical equivalences are identified by natural
  isomorphism, while group, monoid, quotient, splitting, prime, and local
  currency maps remain pointwise data.  Identity, composition, unit laws,
  associativity, representative existence, and exact nonemptiness criteria
  are proved.  Constructing and populating these records from the source
  evaluation and reconstruction algorithms remains open.
  A pointed inversion carries an
  actual decomposition subgroup; restriction along its continuous inclusion
  derives torsion detection from the precise equality
  \(M_{\rm TM}^{\times}\cap\operatorname{res}^{-1}(\mathrm{tors})=
  M_{\rm TM}^{\mu}\), and hence derives the splitting.
  Quotienting an actual constant subgroup by its exact torsion
  subgroup gives a precise multiplication parametrization; it is proved
  bijective when every zero-label theta value is constant torsion.  Given a
  zero-to-base action-pair comparison, the zero-to-label comparisons derive a
  graph proved equal to the diagonal range from Corollary~3.5(i).  The zero
  factor is continuously and multiplicatively equivalent to this graph,
  equivariant for all transported labeled actions.
  Multiplication by \((\mathbb Z/\ell)^\times\) descends to
  the exact EtaleTheta sign-label quotient, factors through
  \((\mathbb Z/\ell)^\times/\{\pm1\}\), and the resulting action on all
  nonzero absolute labels is proved transitive.  After sign identification,
  symmetrization retains one
  topological group/monoid action pair at every nonzero absolute label;
  its outer transports satisfy identity and multiplication coherence, and
  action-compatible isomorphisms derive an injective diagonal monoid map.  Its
  actual range is a submonoid continuously and multiplicatively equivalent to
  the base monoid, and this equivalence intertwines the transported action.
  Thus the label set is not replaced by a coarse one-point quotient.  Its diagonal
  restriction-isomorphism synchronization is represented first by the exact
  quotient for the equivalence relation ``differ by one global inner action.''
  More precisely, the full compatible diagram
  \(\Pi_X\leftarrow\Pi_v\leftarrow G_v^{\langle|F_l|\rangle}\), together with
  its equivariant monoid isomorphism, is quotiented by one base-group element
  that acts simultaneously on the group map and monoid map.  This full
  quotient maps to the monoid-only quotient.  Thus the conjugator cannot
  depend on a label, group element, or monoid element.  Construction of the
  diagram and action from the paper data remains open.  The diagonal
  unit submonoid and finite Gaussian monoid are constructed inside the
  heterogeneous direct product of these labeled factors.  The base group acts
  diagonally through the symmetrizing group isomorphisms; constants are
  equivariant, diagonal units are stable, and profile stability induces an
  actual action by automorphisms of the Gaussian monoid.  For a value
  profile, the homogeneous finite Gaussian monoid is the submonoid of
  the direct product generated by diagonal units and natural profile powers;
  a chosen compatible root system gives the rational-root Gaussian monoid,
  with root ambiguity stated as diagonal torsion.  Root systems differing by
  this permitted ambiguity are proved to generate the same infinite Gaussian
  monoid.  The specialized profile consists of the actual elements
  \(q^{j^2}\).  The \(n\)-torsion of \(\mathbb Q/\mathbb Z\) is proved to
  have exactly \(n\) elements and is explicitly equivalent to the
  roots-of-unity group of a torsion-cyclotomic monoid.  A multiplicative
  realization of the \(\mu_{2\ell}\)-action is required to be injective,
  hence acts freely, and derives its exact torsion-unit range and the
  theta-value family from the continuous \(H^1\) orbits; the resulting
  multiplicative orbit set is proved equal to the original action orbit at
  every label.  Any two
  simultaneous choices from it are proved to differ by independent
  coordinatewise \(2\ell\)-torsion.  Its representative family constructs a
  canonical common \(2\ell\)-multiple Gaussian submonoid.  The ambiguity
  vanishes after taking \(2\ell\)-th powers, so every selected profile yields
  exactly this common object.  The source Hodge--Arakelov pilot places every
  \(q^{j^2}\) in its actual continuous-\(H^1\) theta orbit and every
  compatible positive root in the corresponding infinite-theta orbit.  It
  derives both Gaussian monoids, every \(2\ell\)-element orbit cardinality,
  the full \((2\ell)^{\ell^\ast}\) profile count, and zero-label torsion.
  Multiplicative equivalences preserve the exact
  constant-torsion subgroup, the two quotient factors, and splitting
  bijectivity.  The finite and infinite Gaussian carriers are proved to be
  precisely the products
  \(\Psi_{\rm cns}^{\times}\cdot\xi^{\mathbb N}\) and
  \(\Psi_{\rm cns}^{\times}\cdot\xi^{\mathbb Q_{\geq 0}}\).
  A multiplicative finite-order action proves that a zero-label orbit based at
  the identity consists entirely of torsion.  Thus the actual composite
  continuous-\(H^1\) restriction constructs both finite and infinite
  zero-label splittings once torsion detection on the constant subgroup is
  proved.  More generally, a restriction homomorphism that sends theta values to torsion and detects
  torsion on constants is proved to derive the splitting used in
  Corollaries~1.12, 2.6, 2.8 and Proposition~3.1.  Consequently, once that
  restriction datum and the Corollary~2.8 Gaussian restriction equivalence
  are supplied together with the two factor identifications, both Gaussian
  splittings and their compatibility are derived without accepting a
  Proposition~3.1 splitting proof as input.
  Construction of the paper's
  subgroup/subquotient and decomposition-group homomorphisms, their actions
  and input orbits from the mono-theta system, the symmetrizing isomorphisms,
  root choices, independent conjugacy compatibility, and both restriction
  maps with their torsion and factor-identification proofs remain source
  obligations.
\end{definition}

\begin{definition}[Toy Hodge--Arakelov degree pilots]
  \label{def:hodge-arakelov-concrete-pilots}
  \uses{def:initial-theta-data,def:iut-i-theater-realization}
  \lean{Iut.Stage1.ToyIUTIIHodgeArakelovEvaluationData}
  \lean{Iut.Stage1.ToyIUTIIHodgeArakelovEvaluationData.ofInitialThetaData}
  \lean{Iut.Stage1.ToyIUTIIHodgeArakelovEvaluationData.thetaValueEvaluationSource}
  \lean{Iut.Stage1.ToyIUTIIHodgeArakelovEvaluationData.signOrbitAverage_eq}
  \lean{Iut.Stage1.ToyIUTIIHodgeArakelovEvaluationData.coordinateAverage_translation_invariant}
  \lean{Iut.Stage1.IUTStage1InitialThetaPrimitiveSourcePacket.ofInitialThetaData}
  \leanok
  This finite degree model computes the positive, \(l\)-prime q-order and
  the square-weight profile expected from \(q^{j^2}\).  At each bad place,
  the calculation uses the theta-root torsor, its base-point coordinate
  equivalence, typed
  \(\mathbb Z/l\mathbb Z\) labels, sign classes, balanced-square weights, and
  coordinate, weighted, and sign-orbit averages.  Positivity, normalization,
  sign invariance, and translation invariance are proved.  It does not
  construct the cohomology classes, restriction algorithm, Gaussian monoids,
  or Gaussian Frobenioids of IUT~II, and therefore discharges no paper-level
  evaluation claim.
\end{definition}

\section{Theorem 3.11 and Step (x)}

\begin{definition}[Source-shaped Theorem 3.11 interfaces]
  \label{def:source-theorem-311-interfaces}
  \lean{Iut.SourceMonoAnalyticTensorPacket}
  \lean{Iut.SourceMonoAnalyticTensorPacket.distinguishedSubmodule}
  \lean{Iut.SourceLGPGaussianLogThetaLattice}
  \lean{Iut.SourceLGPGaussianLogThetaLattice.monoAnalyticProcession}
  \lean{Iut.SourceLGPGaussianLogThetaLattice.verticalCoricityToCommon}
  \lean{Iut.SourceLGPGaussianLogThetaLattice.commonMonoAnalyticProcession}
  \lean{Iut.SourceTheorem311LocalData}
  \lean{Iut.SourceTheorem311Ind1Action.ofProcessionFunctor}
  \lean{Iut.SourceTheorem311Ind2Action}
  \lean{Iut.SourceTheorem311LocalInd2Action.action}
  \lean{Iut.SourceTheorem311Ind2SummandAction.packetAction_on_tensor}
  \lean{Iut.SourceTensorPacketKummerIso}
  \lean{Iut.SourcePacketIntegralConstruction.logShellRegion}
  \lean{Iut.SourcePacketIntegralConstruction.logShellRegion_eq_unitBall}
  \lean{Iut.SourceArchimedeanLogShellDefinition.exponential_principalLog}
  \lean{Iut.SourceArchimedeanLogShellDefinition.principalLog_mem_unitBall}
  \lean{Iut.SourceAutHolomorphicSemiGerm.neighborhoodSystem}
  \lean{Iut.SourceAutHolomorphicSemiGerm.iInter_neighborhood_eq_unitCircle}
  \lean{Iut.SourceAutHolomorphicSemiGerm.exists_neighborhood_subset_of_mem}
  \lean{Iut.SourceAutHolomorphicSemiGerm.multiplication_holomorphicOn_levels}
  \lean{Iut.SourceAutHolomorphicSemiGerm.multiplication_tendsto_unitNeighborhoodFilter}
  \lean{Iut.SourceAutHolomorphicSemiGerm.inversion_holomorphicAtLevel}
  \lean{Iut.SourceAutHolomorphicSemiGerm.inversion_tendsto_unitNeighborhoodFilter}
  \lean{Iut.SourceAutHolomorphicSemiGerm.puncturedNeighborhood_eq_inner_union_outer}
  \lean{Iut.SourceAutHolomorphicSemiGerm.innerSide_eq_connectedComponentIn}
  \lean{Iut.SourceAutHolomorphicSemiGerm.outerSide_eq_connectedComponentIn}
  \lean{Iut.SourceAutHolomorphic.autHolomorphicSubgroup}
  \lean{Iut.SourceAutHolomorphic.AutHolomorphic.restrict_mul}
  \lean{Iut.SourceAutHolomorphic.AutHolomorphic.restrict_inv}
  \lean{Iut.SourceAutHolomorphicSemiGerm.levelAutHolomorphic}
  \lean{Iut.SourceAutHolomorphicSemiGerm.levelInclusion_hasAmbientHolomorphicLift}
  \lean{Iut.SourceAutHolomorphicRigidity.Automorphism.exists_integer_exponent}
  \lean{Iut.SourceAutHolomorphicRigidity.Automorphism.exponent_eq_one_or_neg_one}
  \lean{Iut.SourceAutHolomorphicRigidity.Automorphism.toAutHolomorphic_eq_identity}
  \lean{Iut.SourceAutHolomorphicRigidity.Automorphism.ambientGerm_eq_identity}
  \lean{Iut.SourceArchimedeanLogShellDefinition.semiGermSelectedSide_eq_connectedComponentIn}
  \lean{Iut.SourceArchimedeanLogShellDefinition.semiGermNeighborhoodSystem_cofinal}
  \lean{Iut.SourceArchimedeanLogShellDefinition.semiGermMultiplication_tendsto}
  \lean{Iut.SourceArchimedeanLogShellDefinition.semiGermInversion_tendsto}
  \lean{Iut.SourceArchimedeanLogShellDefinition.SemiGermGroupAutomorphism.rigid}
  \lean{Iut.SourceArchimedeanPacketUnitData.summandSemiGermNeighborhoodSystem}
  \lean{Iut.SourceArchimedeanPacketUnitData.UnitTuple.semiGermUnit}
  \lean{Iut.SourceArchimedeanPacketUnitData.principalLogTensor_mem_unitBall}
  \lean{Iut.SourceArchimedeanPacketLogLinkStep}
  \lean{Iut.SourceGaloisLocalUnitData}
  \lean{Iut.SourceNonarchimedeanLogShellDefinition.invariantLocalUnits_fixed}
  \lean{Iut.SourceNonarchimedeanPacketUnitLogData.packetUnitLogSubgroup}
  \lean{Iut.SourceNonarchimedeanPacketUnitLogData.packetUnitLogSubgroup_le_packetLattice}
  \lean{Iut.SourceNonarchimedeanPacketLogLinkStep}
  \lean{Iut.SourceNonarchimedeanPacketLogLinkStep.graph_le_packetUnitLogSubgroup_prod}
  \lean{Iut.SourceNonarchimedeanPacketLogLinkStep.related_of_unitTuplesRelated}
  \lean{Iut.SourceForwardRelation}
  \lean{Iut.SourceForwardDomain}
  \lean{Iut.SourceForwardRange}
  \lean{Iut.sourceForwardIndex_sub_eq}
  \lean{Iut.SourceAbsoluteLGPGaussianLogThetaLattice.verticalPacketUnitLogSubgroup_image}
  \lean{Iut.SourceTheorem311VerticalLocalKummer}
  \lean{Iut.SourceTheorem311VerticalLogLink.logarithm_eq_of_differsByRootOfUnity}
  \lean{Iut.SourceTheorem311VerticalLogLinkFamily.adjacent_logarithm_eq_on_overlap}
  \lean{Iut.SourceTheorem311CoricContainer}
  \lean{Iut.SourceTheorem311Ind3System}
  \lean{Iut.SourceTheorem311Ind3System.nonarchLogIterateRelation}
  \lean{Iut.SourceTheorem311Ind3System.nonarchLogIterateDomain}
  \lean{Iut.SourceTheorem311Ind3System.nonarchLogIterateRange}
  \lean{Iut.SourceTheorem311Ind3System.nonarch_logKummer_contained}
  \lean{Iut.SourceTheorem311Ind3System.nonarch_logKummer_contained_of_mem_range}
  \lean{Iut.SourceTheorem311Ind3System.logShellRegion_eq_unitBall}
  \lean{Iut.SourceTheorem311Ind3System.archLogIterateRelation}
  \lean{Iut.SourceTheorem311Ind3System.archLogIterateDomain}
  \lean{Iut.SourceTheorem311Ind3System.archLogIterateRange}
  \lean{Iut.SourceTheorem311Ind3System.archLogBallPreimage_subset_logShellRegion}
  \lean{Iut.SourceTheorem311Ind3System.archLogBall_relationally_surjective}
  \lean{Iut.SourceTheorem311Ind3System.arch_direct_units_contained}
  \lean{Iut.SourceTheorem311Ind3System.arch_direct_localBalls_contained}
  \lean{Iut.SourceLGPSplittingMonoidAction}
  \lean{Iut.SourceTheorem311BadPrimeLogKummerData}
  \lean{Iut.SourceTheorem311BadPrimeLogKummerData.AdmissibleAcrossLog}
  \lean{Iut.SourceTheorem311BadPrimeLogKummerData.logarithm_eq_of_admissibleAcrossLog}
  \lean{Iut.SourceTopologicalGroupMonoidActionPair.Hom}
  \lean{Iut.SourceKappaLocalizationDiagram}
  \lean{Iut.SourceKappaLocalizationDiagram.inclusion_isEmbedding}
  \lean{Iut.SourceKappaLocalizationDiagram.inclusion_injective}
  \lean{Iut.SourceKappaLocalizationDiagramIso}
  \lean{Iut.SourceKappaLocalizationDiagramIso.localizationSquare}
  \lean{Iut.SourceKappaLocalizationDiagramIso.inclusionSquare}
  \lean{Iut.SourceTheorem311LabeledData}
  \lean{Iut.SourceTheorem311LabeledKummerIso}
  \lean{Iut.SourceTheorem311MultiradialPresentation}
  \lean{Iut.SourceTheorem311MultiradialAlgorithm.Quotient}
  \lean{Iut.SourceTheorem311PermutationBiCoricCompatibility.quotientEquiv}
  \lean{Iut.SourceFPrimeStripHom.NaturallyIsomorphic}
  \lean{Iut.SourceFPrimeStripFullPolyIsomorphism}
  \lean{Iut.SourceFPrimeStripFullPolyIsomorphism.comp_assoc}
  \lean{Iut.SourceFPrimeStripPolyIsomorphismSquare}
  \lean{Iut.SourceFTimesMuPrimeStripFullPolyIsomorphism}
  \lean{Iut.SourceFTimesMuPrimeStripFullPolyIsomorphism.comp_inverse_assoc}
  \lean{Iut.SourceFTimesMuPrimeStripPolyIsomorphismSquare}
  \lean{Iut.SourceFMonoAnalyticPrimeStripFullPolyIsomorphism}
  \lean{Iut.SourceFMonoAnalyticPrimeStripFullPolyIsomorphism.canonical}
  \lean{Iut.SourceFMonoAnalyticPrimeStripFullPolyIsomorphism.canonical_inverse}
  \lean{Iut.SourceFMonoAnalyticPrimeStripFullPolyIsomorphism.comp_assoc}
  \lean{Iut.SourceFTimesMuReconstructionFullness}
  \lean{Iut.SourceFTimesMuReconstructionAlgorithm}
  \lean{Iut.SourceFTimesMuReconstructionAlgorithm.mapFullPolyIsomorphism}
  \lean{Iut.SourceFTimesMuReconstructionAlgorithm.mapFullPolyIsomorphism_comp}
  \lean{Iut.SourceFTimesMuReconstructionAlgorithm.mapFullPolyIsomorphism_surjective}
  \lean{Iut.SourceFTimesMuReconstructionAlgorithm.mapFullPolyIsomorphism_inverse}
  \lean{Iut.SourceFTimesMuPrimeStripPolyIsomorphismSquare.ofUpperHorizontal}
  \lean{Iut.SourceIndexedHorizontalKummerSquare}
  \lean{Iut.SourceTheorem311TrianglePrimeStripSquare}
  \lean{Iut.SourceTheorem311EnvironmentPrimeStripFamily}
  \lean{Iut.SourceTheorem311TimesMuPrimeStripFamily}
  \lean{Iut.SourceTheorem311TimesMuPrimeStripConstruction.toFamily}
  \lean{Iut.SourceTheorem311TimesMuPrimeStripConstruction.thetaLink_full}
  \lean{Iut.SourceTheorem311TimesMuPrimeStripConstruction.toCommon}
  \lean{Iut.SourceTheorem311TimesMuPrimeStripConstruction.fromCommon}
  \lean{Iut.SourceTheorem311TimesMuPrimeStripConstruction.toTriangleSquare}
  \lean{Iut.SourceTheorem311TimesMuTrianglePrimeStripSquare}
  \lean{Iut.SourceTheorem311EnvironmentMonoAnalyticPrimeStripFamily.toTimesMu}
  \lean{Iut.SourceTheorem311EnvironmentTimesMuPrimeStripFamily}
  \lean{Iut.SourceTheorem311EnvironmentTimesMuPrimeStripSquare.ofTriangle}
  \lean{Iut.SourceTheorem311MonoThetaProjectiveSystemFamily}
  \lean{Iut.SourceTheorem311LabeledHorizontalKummerSquare}
  \lean{Iut.SourceTheorem311LabeledHorizontalKummerSquare.localizationAutomorphismEquiv}
  \lean{Iut.SourceTheorem311LabeledHorizontalKummerSquare.inclusionAutomorphismEquiv}
  \lean{Iut.SourceTheorem311HorizontalCorridorBoundary}
  \lean{Iut.SourceTheorem311HorizontalCorridorBoundary.primeStripClauseA}
  \lean{Iut.SourceTheorem311HorizontalEvaluationCompatibility.descend}
  \leanok
  The source layer now represents a local packet by the genuine dependent
  tensor product of its \(j+1\) place-indexed direct sums.  The
  \(I_{\mathbb Q}(S_{j+1,j};v)\) submodule is the span of pure tensors whose
  distinguished factor lies in the summand indexed by \(v\).  Integral
  regions must be either nonarchimedean lattices or archimedean seminorm unit
  balls and must prove their rational spans.  The \(\mathbb Z^2\)-indexed
  lattice contains actual F-theta bridges and distinct Hodge theaters; its
  horizontal theta-links, mono-analytic processions, and vertically coric
  procession maps are constructed.  \(\Ind{1}\) is derived from a functor on
  the category of processions.  \(\Ind{2}\) accepts only the independent
  actions on individual place summands; the direct-sum action, tensor action,
  full dependent-product action, and their laws are derived.  The complete
  presentation also carries its bad-place splitting monoids, number-field
  embeddings, and objects of the four Frobenioid categories.  Its exact
  quotient is the equivalence closure generated by Ind1 and Ind2.  A bi-coric
  bridge member induces a procession isomorphism, and a transport intertwining
  both actions descends, by a proved relation-induction argument, to an
  equivalence of these quotients.

  The vertical interface uses placewise continuous linear Kummer
  equivalences, their induced tensor maps, complete localization diagrams of
  topological group/pseudo-monoid action pairs, number-field ring equivalences,
  splitting-monoid equivalences, and
  structure-preserving equivalences of the realified and non-realified split
  Frobenioids.  \(\Ind{3}\) fixes one vertically coric \(n,\circ\) container
  while \(m\) varies.  At nonarchimedean places it contains both direct Kummer
  images and forward iterated-log images.  The direct finite-place subgroup is
  no longer supplied as arbitrary packet data: it is generated by pure tensors
  of the summandwise invariant-unit logarithm images.  Each such invariant-unit
  type is itself the fixed subgroup of an explicit continuous profinite Galois
  action, rather than an unrelated type carrying an invariant name.  The
  generated packet subgroup lies in the closed tensor lattice and is
  transported exactly by the vertical Kummer equivalence.  A finite log-link
  step is represented as a partial correspondence between the actual fixed
  units in every summand.  Lean applies the summand logarithms, tensors related
  tuples, takes the additive closure of the resulting graph, and recursively
  composes adjacent steps.  It proves that an (r)-fold iterate starting at
  (m-r) ends at (m), and derives containment of the endpoint's Kummer image
  from its unit-log subgroup membership.  Thus neither a direct multi-step map
  nor a final iterated-containment proof is accepted by the concrete Ind3
  constructor.  At an archimedean summand the Hermitian seminorm is normalized
  by \(\pi\), so its unit ball is exactly Definition~1.1(ii)'s raw complex
  radius-\(\pi\) disk; there is no second radius-\(\pi\) ball of the normalized
  metric.  Lean uses \(\arg(z)i\) to construct a principal logarithm in this
  ball for every genuine \(S^1\)-unit and proves that exponentiating recovers
  the unit.  Choosing one place per tensor factor, the direct-sum and tensor
  metric equations put the resulting pure unit tensor, and every selected
  local radius-\(\pi\) tensor, in the packet unit ball.  Each unit is also
  canonically a point of the exact \(S^1\) core of an antitone projective system
  of open annuli.  Lean identifies their intersection with \(S^1\) and proves
  that they are cofinal among all neighborhoods of \(S^1\), restricts ambient
  complex differentiability to every level, proves that ambient multiplication
  and inversion preserve the neighborhood filters, and proves that deleting
  \(S^1\) leaves the two radial
  connected components.  The component determined by
  \(\mathcal O^{\triangleright}_{\mathbb C}\setminus S^1\) is selected by a
  connected-component theorem rather than an arbitrary input.  For every
  connected annulus \(U_n\), Lean now derives the actual group
  \(\operatorname{Aut}_{\mathrm{hol}}(U_n)\) from biholomorphic self-maps and
  proves identity, composition, inversion, and restriction laws.  For a
  groupification-induced holomorphic automorphism of \(\mathbb C^\times\),
  pullback along the exponential covering yields an integer exponent;
  bijectivity restricts it to \(1\) or \(-1\), and preservation of the selected
  inner component rules out inversion.  The induced filter germ is therefore
  proved equal to the identity.  An infinite-place log-link step is only a
  partial correspondence on these semi-germ core-point tuples.
  Recursive composition derives every positive iterate, its domain and range,
  the participating endpoint-ball subset, and relational surjectivity onto the
  source domain.  Direct Kummer containment follows from the full vertical
  unit-ball image.  Thus the concrete Ind3 constructor accepts no arbitrary
  archimedean subgroup, multi-step domain, preimage set, iterate map,
  surjectivity proof, or direct containment.  The general Riemann-surface and
  RC-holomorphic local-morphism category of
  \emph{Absolute Anabelian Topics III}, Definition~2.1 and Corollary~2.3,
  extension of rigidity from groupification-induced maps to every purely local
  group-germ automorphism, and co-holomorphic Kummer reconstruction are not yet
  formalized.  Equality to the common
  measure-derived log-volume implies independence of \(m\).  A multiplicative log-link into
  a rational module provably kills roots of unity.  At a bad prime, cross-link
  the actual unit groups of the splitting monoids map multiplicatively into
  the log-link domain.  Cross-link admissibility is derived from membership in
  the current and next domains together with an explicit root-of-unity quotient
  in the next domain; it is not an arbitrary relation field.  Admissible inputs
  have the same additive log image.
  Horizontal compatibility is no longer an unindexed field of a single
  column.  Its boundary contains two column assemblies, proves that they use
  the same lattice at indices \(n\) and \(n+1\), and fixes all four corners of
  every square.  The active clauses~(iii)(a)--(b) interface now accepts only
  the Definition~4.9 \(F^{\vdash\times\mu}\)-prime-strips.  Each triangle
  corner is computed by one functorial reconstruction algorithm from the
  mono-analytic transport of its actual site or vertically coric theater.  The
  same algorithm computes every environment strip and comparison from the
  Proposition~2.1(vi) mono-analytic data and derives the inverse laws.  Thus no
  prebuilt times-\(\mu\) family enters the active corridor.  The ordinary-\(F\)
  square remains only as precursor scaffolding.  Following the convention of
  IUT~I, Section~0, a
  times-\(\mu\) prime-strip full
  poly-isomorphism is represented by a componentwise
  equivalence map modulo natural isomorphism on its categorical components;
  it is not a strict categorical isomorphism of prime-strip records.  Lean
  proves that identity and composition descend to these classes, together
  with associativity and cancellation of selected inverse pairs.  A
  prime-strip square quantifies in both directions over all horizontal
  classes; its chosen upper and lower classes only witness nonemptiness.  The
  Corollary~4.10(iv) is retained as a representative-level rigidity law for
  every actual adjacent horizontal pair: lifting the underlying mono-analytic map of any output
  equivalence recovers that output up to structured natural isomorphism.  Lean
  proves from this law that every times-\(\mu\) horizontal class has a
  mono-analytic preimage.  The corridor accepts neither a clause~(a) square,
  a chosen theta-link member, an arbitrary quotient-level fullness witness,
  nor chosen Kummer maps.  Every mono-analytic strip carries structure-preserving maps to
  and from the common model with two-sided Section~0 inverse laws.  Lean
  therefore derives the canonical theta-link and both Theorem~1.5 Kummer
  directions, proves their cancellation, reconstructs the upper and vertical maps, defines the
  lower map by Kummer conjugation, and proves commutativity and both
  all-class transport directions.  The triangle
  and environment prime-strip squares range over every \(m\); the environment
  family carries two-sided natural comparison classes with the triangle
  strips, and its square is proved class-by-class in both directions by
  conjugating the triangle square and cancelling those comparisons.  These
  facts validate the exact typed coarsification, reconstruction, fullness
  boundary, and conjugation mechanism.  Constructing the finite/archimedean
  outputs, proving representative-level reconstruction fullness, and deriving
  the coherent model comparisons from Corollaries~4.5, 4.6, 4.10 and
  Theorem~1.5 remain open.  The
  mono-theta square ranges over every \(m\) and bad
  place and uses actual projective systems with levelwise environment
  realizations.  Clause~(iii)(d) no longer collapses
  \(\{\pi_1^{\kappa\text{-sol}}\ltimes M_\infty^\kappa\}\) and its localizations
  to two profinite groups.  For every \(m\), absolute label, and selected place,
  it retains the global action pair, the local \(M_\infty^\kappa\) and
  \(M_\infty^{\kappa\times}\) action pairs, the equivariant localization map,
  and the equivariant topological embedding that realizes the displayed
  inclusion.  These are genuine partial-product
  pseudo-monoids, not total-monoid surrogates.  Horizontal and vertical Kummer
  isomorphisms preserve both arrows, and site/common coherence is equality of
  the complete diagram transports.  The resulting arrow-category squares
  transport the full automorphism groups.  Constructing these diagrams from
  Corollary~4.8 and the cyclotomic-rigidity algorithm remains open.
  Compatible evaluations descend
  to the equivalence relation generated jointly by \(\Ind{1}\)--\(\Ind{3}\).

  These declarations remain a partial construction of Theorem~3.11.  The
  tensor packets, vertical Kummer maps, finite integral images, and finite
  direct unit-log containment are constructed from the typed lattice data.
  Constructing that lattice data, its summand unit/log inputs, splitting
  monoids, fields, and Frobenioids from the preceding source theorems remains
  open, as do the general RC-holomorphic local-morphism classification and
  rigidity for arbitrary purely local group-germ automorphisms,
  construction of the adjacent finite and infinite unit correspondences from
  the Gaussian log-theta lattice, and construction of the indexed horizontal
  isomorphisms and their equivariance from the actual theta links.
\end{definition}

\begin{definition}[Source full vertical log-links]
  \label{def:source-full-vertical-log-links}
  \uses{def:source-theorem-311-interfaces}
  \lean{Iut.SourceSelectedLocalFullLogLink}
  \lean{Iut.SourceSelectedLocalFullLogLink.NonarchimedeanRelated}
  \lean{Iut.SourceSelectedLocalFullLogLink.ArchimedeanRelated}
  \lean{Iut.SourceSelectedLocalFullLogLink.nonarchimedean_target_unique}
  \lean{Iut.SourceSelectedLocalFullLogLink.archimedean_target_unique}
  \lean{Iut.SourceAbsoluteLGPVerticalFullLogLinks}
  \lean{Iut.SourceAbsoluteLGPVerticalFullLogLinks.toVerticalUpperSemiData}
  \leanok
  A selected-place full log-link in the sense of IUT~III,
  Definition~1.1, consists of a continuous rational-linear post-composition
  between the actual ind-topological logarithmic modules, a compatible ring
  equivalence of their realized logarithmic fields, and continuous injective
  realizations of the adjacent target's local units.  At a finite place the
  adjacent relation asserts that the source \(p\)-adic logarithm equals a
  realized target invariant unit.  At an infinite place it asserts the same
  equality for the source principal logarithm and a realized target circle
  unit.  These are partial relations: existence is not assumed, while
  injectivity proves uniqueness of any target.

  A family of these links on all adjacent lattice sites derives the finite
  packet graph and archimedean semi-germ relation componentwise, and hence
  constructs the complete \(Ind{3}\) upper-semi input.  Constructing the
  full links from the cited \(Psi_{\mathrm{cns}}\), Frobenioid, and
  poly-isomorphism data remains a source obligation.
\end{definition}

\begin{definition}[Assembled source Theorem 3.11 column and corridor boundaries]
  \label{def:source-theorem-311-column-boundary}
  \uses{def:source-theorem-311-interfaces,def:source-full-vertical-log-links}
  \lean{Iut.SourceTheorem311ColumnBoundary}
  \lean{Iut.SourceTheorem311ColumnBoundary.presentation}
  \lean{Iut.SourceTheorem311ColumnBoundary.ind1Action}
  \lean{Iut.SourceTheorem311ColumnBoundary.ind2Action}
  \lean{Iut.SourceTheorem311ColumnBoundary.ind3System}
  \lean{Iut.SourceTheorem311ColumnBoundary.verticalLocalKummer}
  \lean{Iut.SourceTheorem311ColumnBoundary.badPrime_logarithm_eq}
  \lean{Iut.SourceTheorem311HorizontalCorridorBoundary}
  \lean{Iut.SourceTheorem311HorizontalCorridorBoundary.siteKappaClauseD_localization_commutes}
  \lean{Iut.SourceTheorem311HorizontalCorridorBoundary.siteKappaClauseD_inclusion_commutes}
  \lean{Iut.SourceTheorem311HorizontalCorridorBoundary.commonKappaClauseD_localization_commutes}
  \lean{Iut.SourceTheorem311HorizontalCorridorBoundary.commonKappaClauseD_inclusion_commutes}
  \leanok
  At one fixed horizontal column, the source assembly forces the
  multiradial presentation and the vertical family to use the same common
  local construction.  Its inputs are the absolute-label Gaussian
  log-theta lattice, mono-analytic log-shell algorithm, finite-stage integral
  data, full local log-links, splitting and labeled Kummer data, and the
  multiradial functor.  From these Lean derives
  the adjacent finite and archimedean packet relations,
  the actual procession action, independent summand action, generated
  \(\Ind{1}/\Ind{2}\) quotient, vertical packet Kummer maps, \(\Ind{3}\)
  system, vertical log-volume independence, and roots-only bad-prime logarithm
  equality.  A separate corridor joins two such boundaries on the same
  lattice at adjacent columns and carries the universally indexed horizontal
  squares and their coherence with vertical Kummer.  Its labeled clause~(d)
  data preserve the full global/local Corollary~4.8 localization diagram and
  every local inclusion, rather than a collapsed profinite-group arrow.

  This assembly is not an existence proof for its upstream inputs.  Its role
  is to replace the unrelated finite adapters by one coherent source boundary
  and to identify precisely which constructions from IUT~I--III remain.
\end{definition}

\begin{definition}[Indeterminacy generator interface]
  \label{def:indeterminacy-generator-interface}
  \lean{Iut.Stage1.IUTStage1Theorem311IndeterminacySourceData}
  \leanok
  The Stage 1 interface names the three source-facing indeterminacy
  generators separately: procession automorphism, local tensor symmetry, and
  upper-semi compatibility.  This is a Lean-checked naming interface, not yet
  a construction of the source Step (x) estimates.
\end{definition}

\begin{definition}[Step (x) averaged corridor record]
  \label{def:step-x-averaged-corridor-record}
  \lean{Iut.Stage1.IUTStage1ProcessionNormalizedIndeterminacyCorridor,Iut.Stage1.IUTStage1ProcessionNormalizedIndeterminacyCorridor.before_average_le_ind3UpperBound}
  \leanok
  The averaged corridor record is the conditional Lean shape of Step (x):
  \(\Ind{1}\) and \(\Ind{2}\) preserve the procession-normalized average, and
  \(\Ind{3}\) supplies the upper-semi bound.  Its Lean theorem proves the
  resulting averaged inequality from those record fields.
\end{definition}

\begin{proposition}[Step (x) corridor]
  \label{prop:step-x-corridor}
  \uses{def:indeterminacy-generator-interface,def:step-x-averaged-corridor-record}
  \notready
  The source Step (x) data should produce a corridor
  \[
    A_0 = A_1 = A_2 \le U_{\Ind{3}},
  \]
  where \(A_0,A_1,A_2\) are the averaged log-volumes before
  indeterminacy, after \(\Ind{1}\), and after \(\Ind{2}\), and
  \(U_{\Ind{3}}\) is the upper-semi target coming from \(\Ind{3}\).
\end{proposition}

\begin{proof}
  In the current code this is an interface-level object.  The proof from the
  source papers must still account for the log-volume invariance under
  \(\Ind{1}\) and \(\Ind{2}\), the upper-semi orientation of \(\Ind{3}\), and
  the normalization conventions for averaging over labels.
\end{proof}

\begin{definition}[Legacy finite M3 procession model]
  \label{def:constructed-m3-procession}
  \uses{def:hodge-arakelov-concrete-pilots,prop:step-x-corridor}
  \lean{Iut.Stage1.IUTIIIProcessionChoice}
  \lean{Iut.Stage1.IUTIIIProcessionChoice.typedCore}
  \lean{Iut.Stage1.IUTIIIProcessionChoice.quotientPossibleImage}
  \lean{Iut.Stage1.IUTIIIVerticalLogKummer.correspondenceLaws}
  \lean{Iut.Stage1.IUTIIIHorizontalCompatibility.ofM2}
  \leanok
  The numerical Gaussian pilots construct a finite coordinate model with two
  actual \(\mathbb Z/l\mathbb Z\)-translation actions.  Their generated
  equality quotient preserves orbit averages and nonempty possible images.
  Natural-number lifts give the correctly oriented \(\Ind{3}\) upper-semi
  relation and do not collapse in that quotient. At every model point,
  the bad-local quotient torsor maps bijectively to the evaluated vertical
  target; its graph, range, and log-volume image are exact, and its finite
  discrete image is compact-open.  The horizontal object retains the distinct
  zero/one histories and carries the typed theta link, gluing, Gaussian
  evaluation, and finite transport laws named \(\IPL/\SHE/\APT\). This
  definition is a regression model and does not discharge Theorem~3.11.
\end{definition}

\begin{theorem}[Legacy degree-level Theorem 3.11 wrapper]
  \label{thm:theorem-311-source-output}
  \uses{def:initial-theta-data,def:hodge-theater-histories,def:prime-strip-links,def:constructed-m3-procession}
  \lean{Iut.Stage1.IUTIIITheorem311ConstructedSource}
  \lean{Iut.Stage1.IUTIIITheorem311ConstructedSource.ofInitialThetaData}
  \lean{Iut.Stage1.IUTStage1InitialThetaPrimitiveSourcePacket.ofInitialThetaData}
  \leanok
  Initial \(\Theta\)-data constructs a finite degree-level regression wrapper
  shaped after Theorem 3.11, including model multiradial possibilities and
  \(\Ind{1}\), \(\Ind{2}\), and \(\Ind{3}\) behavior.  The older qualitative
  packet is not a field of this source and is available only through an
  explicitly named legacy adapter.  It does not construct the source objects
  of Definition~\ref{def:source-theorem-311-interfaces} and does not prove
  Theorem~3.11.  Analytic Step (xi) hull and determinant enrichment is outside
  this wrapper.
\end{theorem}

\section{Step (xi)}

\begin{definition}[Simultaneous holomorphic expressibility]
  \label{def:she}
  \notready
  The \(\SHE\) property says that the relevant output possibilities may be
  expressed relative to the arithmetic holomorphic structure in the \(1\)-column
  while keeping the \qPilot{} log-volume fixed as input data.
\end{definition}

\begin{definition}[Algorithmic parallel transport]
  \label{def:apt}
  \notready
  The \(\APT\) property records how the multiradial algorithm transports data
  through gluings up to indeterminacy.  It belongs primarily to the implication
  from Theorem 3.11 to the reorganized hull-plus-determinant checkpoint.
\end{definition}

\begin{definition}[Hull plus determinant checkpoint]
  \label{def:hdd}
  \notready
  The \(\HDD\) checkpoint is the composite of descent with holomorphic hull
  and determinant formation.  It is the source-side operation that turns
  multiradial output regions into objects comparable with the \qPilot{} at
  the level of signed real log-volume.
\end{definition}

\begin{proposition}[Step (xi) upper-ray input]
  \label{prop:step-xi-upper-ray-input}
  \uses{thm:theorem-311-source-output,def:she,def:hdd}
  \notready
  The multiradial construction, followed by holomorphic hull, determinant, and
  log-volume, should produce a signed upper ray
  \[
    \R_{\le -|\log(\Theta)|}
  \]
  of possible output log-volumes, and the \qPilot{} reading
  \(-|\log(q)|\) should lie in this ray.
\end{proposition}

\begin{proof}
  This is the central source obligation.  It is exactly where the comparison
  must avoid identifying the arithmetic holomorphic histories too early while
  still obtaining a common log-volume container in which the final inequality
  is meaningful.
\end{proof}

\begin{proposition}[Localized Step (xi) \(C_\Theta\) source]
  \label{prop:localized-step-xi-ctheta-source}
  \uses{prop:step-xi-upper-ray-input}
  \lean{Iut.Stage1.Experiments.concreteHodgeTheaterLogThetaPrincipalValuationBallThetaClassFiberTransportThetaRegionDefinedPrincipalValuationBallCompactOpenLogKummerMapPossibleRegionFamilyExactXiSelectedQRemark395SourceDerivedCThetaAuditViaGraphPointwiseLogShellExactCompactOpenImage}
  \lean{Iut.Stage1.Experiments.concreteHodgeTheaterLogThetaPrincipalValuationBallThetaClassFiberTransportThetaRegionDefinedPrincipalValuationBallCompactOpenLogKummerMapPossibleRegionFamilyExactXiSelectedQRemark395ValuationBallHaarValuationCoverFiberThetaRegionLocalArithmeticStepXILocalTermClosedCorollary312Endpoint}
  \leanok
  The current source-derived exact-\(\Xi(P_B)\) route starts from compact-open
  log-Kummer map data over the transported theta-region family and constructs
  the graph, pointwise log-shell, valuation-ball, and Remark~3.9.5
  exact-family audit internally.  The strongest local-arithmetic
  valuation-cover boundary then starts from a named localized vector-bundle
  determinant source: its weighted adjusted local terms give the Step~(xi)
  finite sum for the canonical \(C_\Theta\)-scale, while the local arithmetic
  upper term is decomposed into the Step~(xi) term, a Haar normalization
  defect, and a main-log term.  The localized-determinant constructor removes
  the separate determinant-source handoff: the constructed Remark~3.9.5 source
  is projected from the same localized vector-bundle determinant source.
  \lean{Iut.Stage1.IUTStage1SourcePackage.IUTStage1Remark395ConstructedHolomorphicHullDeterminantSource.ofLocalizedVectorBundleDeterminantSource}
  \lean{Iut.Stage1.IUTStage1SourcePackage.IUTStage1Remark395ConstructedHolomorphicHullDeterminantSource.ofLocalizedVectorBundleDeterminantSource_endpoint}
  \lean{Iut.Stage1.Experiments.concreteHodgeTheaterLogThetaQuotientThetaPilotSource_toOneSidedLocalizedDeterminantStepXILocalTermCThetaSource}
\end{proposition}

\begin{proof}
  Lean first carries the compact-open log-Kummer source through the
  source-derived exact-\(\Xi(P_B)\) global \(C_\Theta\) audit.  At the
  local-term boundary it identifies the localized determinant source with the
  formula-matching Step~(xi) source by construction, derives the displayed
  Ob3/Ob4 and valuation-cover handoffs, and invokes the finite-sum closed
  Corollary~3.12 endpoint.
\end{proof}
